\chapter*{\chapterstyle{II --- Corps des Complexes}} % 99% Fini
\addcontentsline{toc}{section}{Corps des Complexes}

On définit le nombre imaginaire \(i\) dont le carré vaut \(-1\), et on construit alors \(\C\) comme l'extension du corps\footnote[1]{La motivation principale de l'introduction de \(i\) et de cette construction est que \(\C\) est algébriquement clos, ie tout les polynomes de degré \(n\) de \(\C[X]\) ont \(n\) racines.} \(\R\) avec les deux lois usuelles, ie on définit:
\[
   \C := \R[i] = \Bigl\{a + ib \; ; \; a, b \in \R \Bigl\}
\]
On peut alors montrer que c'est un ensemble stable pour les lois usuelles et qu'il vérifie toutes les propriétés qui font de lui un \textbf{corps}.\+
Chaque nombre complexe se définit alors comme des sommes ou produits de réels et du nombre imaginaire et on appelle alors cette expression la \textbf{forme algébrique} d'un nombre complexe et on appelle \(a\) \textbf{la partie réelle} et \(b\) \textbf{la partie imaginaire} de ce nombre. \<

Géométriquement, on peut identifier les nombres complexes à des points du plan, en effet, \(a + ib\) peut se comprendre comme une combinaison linéaire d'un nombre de l'axe réel, et d'un nombre de l'axe imaginaire.

\subsection*{\subsecstyle{Module {:}}}
On appelle \textbf{module} de \(z \in \C\) le \textbf{prolongement} de la fonction valeur absolue à \(\C\), c'est donc une \textbf{norme} et on la définit telle que {:}
\[
   |z| = \sqrt{a^2 + b^2} = \sqrt{z\overline{z}} 
\]
Dans la suite, on notera \(\rho\) le module de \(z\) pour faciliter la lecture.

\subsection*{\subsecstyle{Forme trigonométrique {:}}}
L'interprétation géométrique permet alors de montrer par passage en coordonées polaires qu'il existe un unique angle \(\theta\) (modulo \(2\pi\)) qu'on appelle \textbf{argument} de \(z\) tel que:
\[
    z = \rho(\cos\theta+i\sin\theta)   
\]

\subsection*{\subsecstyle{Forme exponentielle {:}}}
De même on définit alors \textbf{la forme exponentielle} de \(z\) l'expression:
\[
    z = \rho e^{i\theta} := \rho(\cos\theta+i\sin\theta)
\]
On peut alors étendre les propriétés usuelles de l'exponentielle à \(\C\) et on en déduit:
\customBox{width=6cm}{
   \begin{align*}
      \arg(zz') &\eqmod{2\pi} \arg(z) + \arg(z')\\
      \arg(\frac{z}{z'}) &\eqmod{2\pi} \arg(z) - \arg(z')
   \end{align*}
}

\subsection*{\subsecstyle{Conjugué {:}}}
On appelle conjugaison \textbf{l'involution} qui à \(z\) associe son \textbf{conjugué}, noté \(\overline{z}\) tel que:
\[
    \overline{z} := a - bi = \rho(\cos\theta - i\sin\theta) = \rho e^{-i\theta}
\]
C'est une application \textbf{additive} et \textbf{multiplicative}, on montre alors les formules suivantes {:}
\customBox{width=7.5cm}{
   \[
      \Re(z) := \frac{z + \overline{z}}{2}
      \quad\quad\quad\quad\quad\quad
      \Im(z) := \frac{z - \overline{z}}{2i}
   \]
}

En utilisant ces formules pour \(z\) sous forme exponentielle, on a alors les \textbf{formules d'Euler} qui sont très importantes car elle permettent de \textbf{linéariser} des expression trigonométriques.\<

\subsection*{\subsecstyle{Formule de Moivre {:}}}
Un propriété importante des formes trigonométriques et exponentielles apellée \textbf{formule de Moivre}\footnote[1]{Ici, on a choisi de considérer \(z \in \U\) mais ces propriétés sont vraies pour \textbf{tout nombre complexe}, il suffit alors d'appliquer la puissance au module.} est:
\customBox{width=9cm}{
   \begin{align*}
      (e^{i\theta})^n &= e^{n(i\theta)}\\
      (\cos\theta + i\sin\theta)^n &= \cos n\theta + i\sin n\theta
  \end{align*}  
}
\begin{center}
    \textit{Les différentes puissances d'un nombre complexe (de module 1) s'interprétent alors comme des points situés à equidistance sur un cercle.}
\end{center}
Graphiquement:
\begin{center}
   \begin{tikzpicture}[yscale=2, xscale=2]
      \coordinate (O) at (0,0);
      \coordinate (X) at (1.5,0);
      \coordinate (Z1) at (0.86, 0.5);
      \coordinate (Z2) at (0.5, 0.86);
   
      \draw[-latex] (-1.5,0) -- (X) node [right] {$\mathbb{R}$};
      \draw[-latex] (0, -1.5) -- (0,1.5) node [above] {$i\mathbb{R}$};
   
      \draw[color = DarkBlue1, thick] (0,0) circle (1cm);
   
      \draw node[color = BrightRed1] at (1.25, 0.6) {$z = e^{i\cdot\theta_1}$}; 
      \draw node[color = BrightRed1] at (0.95, 1) {$z^2 = e^{i\cdot 2\theta_1}$}; 
   
      \draw[color = BrightRed1!75, thick] (O) -- (Z1);
      \pic [color = BrightRed1, draw, thick, "$\theta_1$", angle eccentricity=1.5, angle radius=0.5cm] {angle = X--O--Z1};
   
      \draw[color = BrightRed1!75, thick] (O) -- (Z2);
      \pic [color = BrightRed1, draw, thick, "$\theta_1$", angle eccentricity=1.5, angle radius=0.75cm] {angle = Z1--O--Z2};
   
      \fill[color = BrightRed1!100] (0.21, 0.67) circle[radius=0.5pt];
      \fill[color = BrightRed1!80] (0.073, 0.70) circle[radius=0.5pt];
      \fill[color = BrightRed1!60] (-0.07, 0.70) circle[radius=0.5pt];
      \fill[color = BrightRed1!40] (-0.21, 0.67) circle[radius=0.5pt];
      \fill[color = BrightRed1!20] (-0.35, 0.61) circle[radius=0.5pt];
   \end{tikzpicture}  
\end{center}

\subsection*{\subsecstyle{Racines n-ièmes {:}}}
Soit \(n \in \N\), une partie importante des problèmes impliquant des nombres complexes proviennent d'équations d'inconnue \(Z\) de la forme:
\[
   Z^n = z
\]
On peut montrer que l'ensemble des solutions de ce type de problème est:
\customBox{width=8cm}{
   \begin{align*}
      S = \Bigl\{ \sqrt[n]{\rho}e^{i\frac{\theta + 2k\pi}{n}}\; ; \; k \in \bigl\{0, 1, \ldots , n-1 \bigl\} \Bigl\}   
  \end{align*}  
}
\underline{Cas particulier {:}}
Si on a une racine n-ième \(Z_0\) de \(Z\) et qu'on connaît les racines n-ièmes de l'unité, alors on peut obtenir toutes les racines n-ièmes de \(Z\) grâce à:

\[
   \Bigl\{ Z \in \C \; ; \; Z^n = z \Bigl\} = \Bigl\{ Z_0u \; ; \; u \in \U_n \Bigl\}   
\]

\subsection*{\subsecstyle{Le nombre complexe \(j\) {:}}}
On note \(j\) la première racine troisième de l'unité.
Le nombre \(j\) est singulier, car il vérifie:
\customBox{width=4cm}{
   \(
      j^2 = j^{-1} = \overline{j}
   \)
}

Graphiquement, on peut observer que les affixes des nombres \(1\), \(j\) et \(\overline{j}\) forment un triangle équilatéral inscrit dans le cercle trigonométrique.

\chapter*{\chapterstyle{II --- Anneau des Polynômes}} % 75% Fini
\addcontentsline{toc}{section}{Anneau des Polynômes}

\subsection*{\subsecstyle{Définition {:}}}

Soit \(K\) un corps, et \(n \in \N\), on appelle \textbf{polynômes} à coefficients dans \(K\) en l'indéterminée \(X\) les éléments de l'ensemble:
\[
   K[X] := \Biggl\{ \; \sum_{i=0}^n a_{i}X^{i} \; ; \; a_i \in K \;\Biggl\}
\]

\subsection*{\subsecstyle{Degré et Valuation {:}}}
Soit \(P, Q, R \in K[X]\).\+
On définit tout d'abord une propriété fondamentale appelée \textbf{degré} de \(P\) telle que deg\((P)\) est le plus grand coefficient non nul de P.
On a alors les propriétés du degré ci-dessous:
\customBox{width=7cm}{
   \begin{align*}
      \deg(P + Q) &\leq \max(\deg(P), \deg(Q))\\
      \deg(PQ) &= \deg(P) + \deg(Q)
  \end{align*}
}

La valuation est définie de manière analogue comme le plus petit coefficient non nul de P.
\subsection*{\subsecstyle{Opérations {:}}}
On considère ci-dessous que deg\((P)\) = \(n\) et deg\((Q)\) = \(m\) et on note \(a_i, b_i\) les coefficients de \(P\) (resp. \(Q\)).\<

On adjoint à cet ensemble une loi d'addition qui est simplement effectuée terme à termes.\+
On adjoint à cet ensemble une loi de multiplication se définit explicitement comme suit:
\[
   PQ := \sum_{k=0}^{n+m}\sum_{i+j = k} a_{i}b_{i}X^{k}
\]

Muni de ces deux opérations, on donne à cet ensemble une structure \textbf{d'anneau}\footnote[1]{Intègre car les coefficients viennent d'un corps.}.\<

On ajoute aussi une opération appelée \textbf{dérivation formelle} d'un polynôme qui s'effectue comme la dérivation analytique usuelle.
\subsection*{\subsecstyle{Divisibilité {:}}}
On définit tout d'abord une \textbf{division euclidienne} de deux polynômes qui se comporte comme la division euclidienne usuelle, à la différence que la condition d'arrêt porte sur le \textbf{degré du reste} qui doit être inférieur à celui du diviseur.\<

Soit \(U, V \in R[X]^2\), on peut aussi définir une relation de \textbf{divisibilité} entre deux polynômes, cette relation est \textbf{transitive et réflexive} et on a aussi:
\customBox{width=6.5cm}{
  \(
      D|A \land D|B \implies D| UA + VB   
  \)
}

\subsection*{\subsecstyle{Racines et Factorisation {:}}}
On dit que \(\alpha\) est une \textbf{racine} de \(P\) si \(P(\alpha) = 0\).\+
On montre alors le théorème fondamental ci-dessous:
\customBox{width=5cm}{
  \(
      P(\alpha) = 0 \Longleftrightarrow (X - \alpha) | P  
   \)
}
\pagebreak

On appelle \textbf{multiplicité d'une racine} \(\alpha\) l'entier \(m\) tel que:
\[
   \Bigr[ (X - \alpha)^m | P \Bigr] \; \land \; \Bigr[(X - \alpha)^{m+1} \nmid P\Bigr]
\]
Si on note \(P^{m}\) la dérivée n-ième de \(P\), on a aussi:
\[
   P^{m}(\alpha) = 0 \; \land \; P^{m + 1}(\alpha) \neq 0
\]

On en déduit que pour une racine \(\alpha\) de multiplicité \(m\), on peut \textbf{factoriser} \(P\) par \((X-\alpha)^m\).\<

Si on considère maintenant plusieurs racines \textbf{distinctes} \(a_0, a_1, \ldots, a_{n-1}, a_n\) de multiplicité respectivement \(m_0, m_1, \ldots, m_{n-1}, m_n\), on peut montrer la propriété suivante:
\begin{flalign*}
   \Bigr[ \prod_{i=0}^{n}(X-\alpha_i)^{m_i} \Bigr] \; \Biggr| \; P \shorteqnote{(On peut factoriser par le produit des \((X-\alpha_i)^{m_i}\))}
\end{flalign*}
\subsection*{\subsecstyle{Décomposition {:}}}
On appelle décomposition d'un polynôme \(P\) une factorisation \textbf{irréductible} de \(P\), cette décomposition dépend du corps considéré, en effet si on considère \(\C[X]\), on peut montrer le \textbf{théorème fondamental de l'Algèbre} ci-dessous:
\customBox{width=11cm}{
   Tout polynôme de degré \(n\) admet \(n\) racines dans \(\C\).
}

Par suite, on peut montrer que tout les polynômes de \(\C[X]\) sont \textbf{scindés}\footnote[1]{C'est à dire que tout les polynômes irréductibles de \(\C[X]\) sont de degré 1.}.\+
Par contre dans \(\R[X]\), il existe évidemment des polynômes de degré 2 irréductibles.\<

Soit \(z \in \C\), il existe une propriété très utile pour décomposer un polynôme \textbf{à coefficients réel} qui est:
\[
   P(z) = 0 \implies P(\overline{z}) = 0
\]
\subsection*{\subsecstyle{Fonctions symétriques des racines {:}}}
On définit le \(k\)-ième \textbf{polynôme symétrique} à \(n\) indéterminées comme la somme des produits à \(k\) facteurs de ses indéterminées, et on le note \(\sigma_k\). \<

Par exemple pour un polynôme en trois indéterminées \(a, b, c\), on a successivement:
\begin{flalign*}
    &\sigma_1 = a + b  + c\shorteqnote{(Somme des produits à 1 facteur)}\\
    &\sigma_2 = ab + ac + bc \shorteqnote{(Somme des produits à 2 facteurs)}\\
    &\sigma_3 = abc\shorteqnote{(Somme des produits à 3 facteurs)}
\end{flalign*}
De manière générale on a:
\[
    \sigma_k(x_1, \ldots, x_n) = \sum_{1 \leq i_1 < \ldots < i_k \leq n} x_{i_1}x_{i_2} \ldots x_{i_k}
\]
Soit \(P \in K[X]\) un polynôme de degré \(n\), on note \(x_1, x_2, \ldots, x_{n-1}, x_n\) ses \(n\) racines et \(a_1, a_2, \ldots, a_{n-1}, a_n\) ses \(n\) coefficients.\+
Alors, pour tout \(k \in \inticc{0}{n}\), on a le \textbf{théorème}:
\customBox{width=8cm}{
   \[\sigma_k(x_1, x_2, \ldots, x_{n-1}, x_n) = (-1)^{k}\frac{a_{n-k}}{a_n}\]
}

Ce théorème permet d'obtenir une relation \textbf{coefficient-racines}.\<

\underline{Exemple:} \(\begin{cases}
    &\text{Pour \(k = 1\), on trouve que la somme des racines de \(P\) vaut \(-\frac{a_{n-1}}{a_n}\)}\\
    &\text{Pour \(k = 2\), on trouve que la somme des doubles produits des racines de \(P\) vaut \(\frac{a_{n-2}}{a_n}\)}\\
    &\ldots\ldots\ldots\ldots\ldots\ldots
\end{cases}\)\<

\chapter*{\chapterstyle{II --- Espaces Vectoriels}} % 99% Fini
\addcontentsline{toc}{section}{Espaces Vectoriels}

\subsection*{\subsecstyle{Définition {:}}}
Soit \(E\) un ensemble \textbf{non-vide}, \(K\) un corps commutatif, et \(u \in E\), on dit alors que \((E, +, \cdot)\) est un \textbf{espace vectoriel sur K} si les conditions ci-dessous sont réunies:
\begin{align*}
   &\bullet \;\;\; (E, +) \text{ est un \textbf{groupe abélien}} \\
   &\bullet \;\; \text{ La loi \(\cdot\) est une application de \(K \times E \longrightarrow E\)} \\
   &\bullet \;\; \text{ La loi \(\cdot\) vérifie \textbf{la distributivité mixte et l'associativité mixte}} \\
   &\bullet \;\;\; 1 \cdot u = u
\end{align*}

On appelle alors \textbf{vecteurs} les éléments de \(E\) et \textbf{scalaires} les éléments de \(K\).

\subsection*{\subsecstyle{Sous-espaces vectoriels {:}}}
Soit \(F \subseteq E\), on dit que \(F\) est un \textbf{sous-espace vectoriel} si et seulement si:
\begin{align*}
   &\bullet \;\; F \text{ est non vide.} \\
   &\bullet \;\; F \text{ est stable par somme.} \\
   &\bullet \;\; F \text{ est stable par multiplication externe.}
\end{align*}

Il est utile de noter que \textbf{l'intersection} de deux sous-espaces vectoriels est aussi un sous-espace vectoriel mais que l'union de deux sous-espaces vectoriels n'est \textbf{en général} pas un sous-espace vectoriel.\<

Un cas particulier est celui du \textbf{sous-espace vectoriel engendré} par \(F\) qu'on note:
\[
   \text{Vect(\(F\))} := \Biggl\{ \sum_{i=0}^{n} \lambda_i u_i\; ; \; \lambda_{i} \in K \; , \; u_{i} \in F \Biggl\}
\]
\begin{center}
    \textit{
        C'est \textbf{l'ensemble des combinaisons linéaires} des vecteurs de \(F\).
    }
\end{center}

Soit \((A_n)_{n \in \N}\) une famille de parties de \(E\), alors on a la propriété suivante:
\customBox{width=5.5cm}{
   \[
      \sum_{i \in \N} \text{Vect(\(A_i\))} = \text{Vect\(\Bigl(\bigcup_{i \in \N} A_i\Bigl)\)}
   \]
}

Une propriété fondamentale est que Vect(\(F\)) est un \textbf{opérateur de clôture} par combinaison linéaire, ie c'est une application \textbf{idempotente, croissante et extensive}.

\subsection*{\subsecstyle{Familles libre et génératrices {:}}}
Soit \(\Fam := (u_1, u_2, \ldots, u_{n-1}, u_n)\) une famille de vecteurs de \(E\).\<

On dit que \(\Fam\) est \textbf{génératrice} de E si on a \(E = \text{Vect(\(\Fam\))}\).\+
On dit que \(\Fam\) est \textbf{libre} si tout combinaison linéaire nulle de vecteurs de \(\Fam\) est à coefficient tous nuls.\<

Formellement, \(\Fam\) est libre si et seulement si on a:
\[
    \forall (\lambda_1, \ldots, \lambda_n) \in K^n \; ; \; \Biggr[ \sum_{i=0}^n \lambda_i u_i = 0_E \implies (\lambda_1, \ldots, \lambda_n) = (0, \ldots, 0) \Biggr]  
\]
\pagebreak

Soit \(\Fam := (u_1, \ldots, u_k, \ldots, u_n)\) une famille libre et génératrice.\<

Si un vecteur \(v\) n'appartient pas à \(\text{Vect(\(\Fam\))}\), alors la famille étendue \((u_1, \ldots, u_n, v)\) est toujours libre.\+
Si un vecteur \(u_k\) appartient à \(\text{Vect(\(\Fam\))}\), alors la famille réduite \((u_1, \ldots, u_{k-1}, u_{k+1}, \ldots, u_n)\) est toujours génératrice.

\subsection*{\subsecstyle{Bases {:}}}
On appelle \textbf{base} de \(E\) une famille \textbf{libre et génératrice}, et on a le théorème suivant:
\customBox{width=8cm}{
   Tout espace vectoriel possède une base.
}

Soit \(\mathscr{B} = (e_1, e_2, \ldots, e_{n-1}, e_{n})\) une base de \(E\).\+
On appelle \textbf{coordonnées} de \(u\) dans la base \(\mathscr{B}\) les coefficients de la décomposition de \(u\) dans la base \(\mathscr{B}\) et on note:
\[
    [u]^{\mathscr{B}} = [\lambda_1 e_1 + \lambda_2 e_2 + \ldots + \lambda_{n-1} e_{n-1} + \lambda_{n-1} e_{n-1} ]^{\mathscr{B}} = 
    \begin{pmatrix}
        \lambda_1\\
        \vdots\\
        \lambda_{n}
    \end{pmatrix}
\]

\subsection*{\subsecstyle{Sommes et espaces supplémentaires{:}}}
Soit \(n\) un entier naturel et \((F_n)_{n \in \N}\) une famille de \(n\) sous-espaces vectoriels de \(E\), considérons la somme ensembliste \(S_n\) telle que:
\[
    S_n := \sum_{n\in\N} F_n := \Biggl\{ \sum_{n\in\N} u_n \; ; \; u_n \in F_n\Biggl\}
\]
Soit \((u_1, u_2, \ldots, u_n)\) un n-uplet de vecteurs quelconques du produit cartésien des \(F_n\), alors la somme est dite \textbf{directe} si et seulement si:
\customBox{width=10cm}{
   \(
      u_1 + u_2 + \ldots + u_n = 0_E \implies (u_1, u_2, \ldots, u_n) = (0, 0, \ldots, 0)
  \)
}

Et on la note alors:
\[
    S_n = \bigoplus_{n\in\N} F_n
\]

Soit \(\B_n\) des bases des \(F_n\), soit la famille \(\Fam_n = (\B_1, \B_2, \ldots, \B_n)\), \textbf{ie les bases des \(F_n\) mises ``bout à bout``}.\+
Alors on a alors le théorème suivant:
\customBox{width=7.5cm}{
   \[
      \Fam_n \text{ est une base de } S_n \Longleftrightarrow S_n = \bigoplus_{n\in\N} F_n
  \]
}

Si \textbf{l'espace entier} est engendré par \(S_n\) et que la somme est \textbf{directe}, alors on dira que les \(F_n\) sont \textbf{supplémentaires} dans \(E\).
    
\subsection*{\subsecstyle{Hyperplans {:}}}
Soit \(F\) un sous-espace vectoriel de \(E\).\+
\(F\) est un \textbf{hyperplan} de E si et seulement si il vérifie l'une des conditions équivalentes suivantes:
\begin{align*}
    &\bullet \text{ Il existe \textbf{une droite vectorielle} \(G\) de \(E\) telle que \(F \oplus G = E\).} \\
    &\bullet \text{ \(H\) est défini par une \textbf{équation linéaire homogène} non triviale.}
\end{align*}
Les hyperplans donc les \textbf{sous-espaces propres maximaux} de \(E\).

\begin{center}
    \textit{
        Résoudre un système d'équations homogènes revient donc à chercher \textbf{l'intersection} des hyperplans que représentent ces équations et donc l'ensemble des solutions est lui aussi nécessairement \textbf{un sous-espace vectoriel}.
    }
\end{center}

\chapter*{\chapterstyle{II --- Applications Linéaires}} % Fini 99%
\addcontentsline{toc}{section}{Applications linéaires}

Soit deux \(\K\)-espaces vectoriels \(E\) et \(F\), et \(f : E \longrightarrow F \).

\subsection*{\subsecstyle{Définition{:}}}
\addcontentsline{toc}{subsection}{Définition}

On dit que \(f\) est une \textbf{application linéaire} ou \textbf{morphisme d'espaces vectoriels} si et seulement si pour tout couple de vecteurs \(u, v \in E\) et pour tout scalaire \(\lambda\) elle vérifie: 
\customBox{width=8cm}{
   \[
      \begin{aligned}
         &\bullet \textbf{ Additivité : } f(u + v) = f(u) + f(v)\\
         &\bullet \textbf{ Homogénéité : } f(\lambda u) = \lambda f(u)\\
     \end{aligned}
  \]
}
On note alors \(\mathcal{L}(E, F)\) l'ensemble des applications linéaires de \(E\) vers \(F\).\<

On appelle \textbf{isomorphisme} tout application linéaire bijective.\+
On appelle \textbf{endomorphisme} tout application linéaire d'un ensemble dans lui-même.\+
On appelle \textbf{automorphisme} tout application linéaire bijective d'un ensemble dans lui-même\footnote[1]{On appelle l'ensemble des automorphismes \textbf{le groupe linéaire} de \(E\) qu'on note \(GL(E)\).}.\<

On appelle \textbf{forme linéaire} tout application linéaire qui a pour ensemble d'arrivée \textbf{le corps de base}\footnote[2]{On appelle l'ensemble des formes linéaires \textbf{l'espace dual} de \(E\) qu'on note \(E^*\).}.

\subsection*{\subsecstyle{Propriétés{:}}}
\addcontentsline{toc}{subsection}{Propriétés}

On s'intéresse au propriétés de l'ensemble \(\mathcal{L}(E, F)\), tout d'abord on peut prouver aisément que c'est \textbf{un sous-espace vectoriel} de l'espace des fonctions, ie il est stable par somme et par multiplication exeterne.\+
En outre, \textbf{la composée de deux applications linéaires est linéaire}.\<

Considérons une application linéaire bijective de \(\mathcal{L}(E, F)\), alors:
\begin{center}
   \textbf{Sa réciproque est nécessairement linéaire\footnote[3]{Découle directement des définitions.}}.
\end{center}

\subsection*{\subsecstyle{Image et Noyau{:}}}
\addcontentsline{toc}{subsection}{Image et Noyau}

On appelle \textbf{image} d'une application linéaire l'image telle que définie pour les applications classiques.\+
On appelle \textbf{noyau} d'une application linéaire tout \textbf{les antécédent du vecteur nul}.\+
\[
   \text{Im}(f) := \Bigl\{ f(u) \; ; \; u \in E \Bigl\}
   \quad\quad\quad\quad\quad
   \text{Ker}(f) := \Bigl\{ u \in E \; ; \; f(u) = 0_F \Bigl\}
\]
On voit directement que \(\text{Im}(f)\) est un \textbf{sous-espace vectoriel} de l'ensemble d'arrivée et que \(\text{Ker}(f)\) est un \textbf{sous-espace vectoriel} de l'ensemble de départ.\<

En particulier si \(E = \text{Vect}(e_1, e_2, \ldots, e_n)\) pour une certaine famille \((e_i)_{i\in \N}\), alors on a\footnote[5]{Idem.}:
\customBox{width=11cm}{
   \(
      \text{Im}(f) = f(\text{Vect}(e_1, e_2, \ldots, e_n)) = \text{Vect}(f(e_1), f(e_2), \ldots, f(e_n))
  \)
}

Une propriété fondamentale du noyau\footnote[6]{De manière générale, cette propriété est vraie pour le noyau de tout morphisme de groupe, on peut alors voir la ``taille'' du noyau comme \textbf{une mesure de la non-injectivité} du morphisme.} est la suivante:
\customBox{width=15.5cm}{
   Une application linéaire est injective si et seulement si son noyau est réduit au vecteur nul.
}
\pagebreak

\subsection*{\subsecstyle{Caractérisations par les familles{:}}}
\addcontentsline{toc}{subsection}{Caractérisations par les familles}

Soit \(\mathscr{F} = (e_i)_{i \in \N}\) une famille de \(E\) et un endomorphisme de \(E\), alors \(f\) est entièrement caractérisée par l'image de cette famille, en effet on a:
\begin{itemize}
   \item L'image d'une famille libre est libre \(\Longleftrightarrow f\) est injective.
   \item L'image d'une famille génératice est génératice \(\Longleftrightarrow f\) est surjective.      
\end{itemize}

\subsection*{\subsecstyle{Endomorphismes remarquables{:}}}
\addcontentsline{toc}{subsection}{Endomorphismes remarquables}

On considère ici des endomorphismes remarquables comme:
\begin{itemize}
   \item Les homotéthies: Elles sont caractérisées par \(\varphi_k : u \longmapsto ku\)
   \item Les symétries: Elles sont caractérisées par \(\varphi \circ \varphi = \text{Id}_E\)
   \item Les projecteurs: Elles sont caractérisées par \(\varphi \circ \varphi = \varphi\)
\end{itemize}
Précisons le cas des projecteurs, en effet si on considère deux sous-espaces \(F, G\) supplémentaires dans \(E\), alors chaque élément \(u \in E\) admet une décomposition unique de la forme \(u = v_F + v_G\).\<

Alors le projecteur de direction \(F\) sur \(G\) est exactement l'application \(\varphi\) telle que:
\[
   \varphi(u) = v_G
\]
Graphiquement, pour \(\varphi\) le projecteur de direction \(F\) sur \(G\):
\begin{center}
   \begin{tikzpicture}[yscale=0.98]
      \draw[black!15] (-4,0) -- (4,0);
      \draw[black!15] (0, -2) -- (0,2);

      \draw[color=BrightBlue1, domain=-0.5:0.5, thick] plot (\x,{4*\x}) node[left] {$F$};
      \draw[color=BrightBlue1, domain=-4:4, thick] plot (\x,{0.5*\x}) node[below right] {$G$};

      \draw[-latex, color=BrightRed1, thick] (0, 0) -- (3/14, 6/7) node[left] {$v_F$};
      \draw[-latex, color=BrightRed1, thick] (0, 0) -- (9/7, 9/14) node[below right] {$v_G = \varphi(u)$};

      \draw[color=BrightRed1, dashed] (9/7, 9/14) -- (1.5, 1.5);
      \draw[color=BrightRed1, dashed] (3/14, 6/7) -- (1.5, 1.5);

      \draw[-latex, color=DarkBlue1, thick] (0, 0) -- (1.5,1.5) node[above left] {$u$};
   \end{tikzpicture} 
\end{center}

\subsection*{\subsecstyle{Applications multilinéaires{:}}}
\addcontentsline{toc}{subsection}{Applications multilinéaires}

On peut généraliser le concept de linéarité à celui de \textbf{multilinéarité} ou \textbf{n-linéarité}. On considère alors une application de la forme \(f : E^n \longrightarrow F\), alors on dit que \(f\) est multilinéaire si et seulement si elle est linéaire \textbf{en chacune des variables}, ie si:
\begin{align*}
   f(e_1 + \lambda u, e_2, \ldots, e_n) &= f(e_1, e_2, \ldots, e_n) + \lambda f(u, e_2, \ldots, e_n)  \\
   f(e_1, e_2  + \lambda u, \ldots, e_n) &= f(e_1, e_2, \ldots, e_n) + \lambda f(e_1, u, \ldots, e_n) \\
   & \vdotswithin{=} \\   
   f(e_1, e_2 , \ldots, e_n + \lambda u) &= f(e_1, e_2, \ldots, e_n) + \lambda f(e_1, e_2, \ldots, u)
\end{align*}
Une telle application est dite:
\begin{itemize}
   \item \textbf{Symétrique} si permuter deux variables préserve le résultat.\footnote[1]{\underline{Exemple:} \(f(e_1, e_2) = f(e_2, e_1)\)}
   \item \textbf{Antisymétrique} si permuter deux variables change le signe du résultat.\footnote[2]{\underline{Exemple:} \(f(e_1, e_2) = -f(e_2, e_1)\), en particulier, le signe du résultat aprés une permutation \(\sigma\) dépends alors de \textbf{la signature de la permutation} (la parité du nombre de permutations effectuées).}
   \item  \textbf{Altérnée} si elle s'annule à chaque fois qu'on l'évalue sur un k-uplet contenant deux vecteurs identiques.\footnote[3]{\underline{Exemple:} \(f(e_1, e_1) = 0\)}
\end{itemize}

\chapter*{\chapterstyle{II --- Théorie de la dimension}} % Fini 99%
\addcontentsline{toc}{section}{Théorie de la dimension}

Dans ce chapitre, on considérera un espace vectoriel \(E\) qui admet une famille génératrice \textbf{finie}. On dira alors que \(E\) est de \textbf{dimension finie}.

\subsection*{\subsecstyle{Théorème de la base incomplète{:}}}
\addcontentsline{toc}{subsection}{Théorème de la base incomplète}

Soit \(\mathcal{L}\) une famille libre et \(\mathcal{G}\) une famille génératrice de \(E\), le concept de \textbf{dimension} se définit gràce au \textbf{théorème de la base incomplète}:
\begin{itemize}
   \item On peut \textbf{compléter} \(\mathcal{L}\) en une base de \(E\) par ajouts de vecteurs de \(\mathcal{G}\).
   \item On peut \textbf{extraire} de \(\mathcal{G}\) une base de \(E\) (corollaire immédiat).
\end{itemize}
Ce théorème permet alors d'assurer \textbf{l'existence} d'une base d'une espace vectoriel de dimension finie.

\subsection*{\subsecstyle{Théorème de la dimension{:}}}
\addcontentsline{toc}{subsection}{Théorème de la dimension}

Il est évident que le cardinal d'une partie libre est toujours inférieur au cardinal d'une partie génératrice \footnote[1]{Aussi appellé \textbf{lemme de Steiniz}.}, de cette considération, on peut alors montrer \textbf{le théorème de la dimension}:
\customBox{width=13cm}{
   Toutes les bases d'un espace vectoriel de dimensions fini ont \textbf{même cardinal}.
}
Ce théorème permet alors d'assurer \textbf{l'unicité} du cardinal des bases d'une espace vectoriel de dimension finie.

\subsection*{\subsecstyle{Définition de la dimension{:}}}
\addcontentsline{toc}{subsection}{Définition de la dimension}

Des deux théorèmes précédents, on a alors l'existence de bases d'une espace de dimension fini, et l'unicité de leur cardinal, on peut alors définir \textbf{la dimension d'une espace vectoriel} comme ce cardinal et on la note \(\text{dim}(E)\).\<

On peut alors en déduire la dimension des espaces vectoriels communs:
\begin{itemize}
   \item L'espace vectoriel classique \(\K^n\) est de dimension \(n\).
   \item L'espace vectoriel des polynômes de degré au plus \(n\) est de dimension \(n + 1\).
   \item Une droite vectorielle est de dimension \(1\).
   \item Un plan vectoriel est de dimension \(2\).
\end{itemize}

\subsection*{\subsecstyle{Espaces de dimension finie{:}}}
\addcontentsline{toc}{subsection}{Espaces de dimension finie}

Considérons maintenant \(E\) un espace vectoriel de dimension finie \(n\) est \(\Fam = (e_i)_{i \in \N}\) une famille de \(n\) vecteurs de \(E\).\+
Alors par déduction immédiate\footnote[2]{En effet par exemple sachant qu'elle est libre et possède \(n\) vecteurs, si ce n'était pas une base, on pourrait rajouter un/des vecteur pour obtenir une base, et donc dim\((E) \neq n \implies \bot\).} de la définition de dimension, on a:
\customBox{width=11cm}{
   \(\Fam\) est libre \(\Longleftrightarrow\) \(\Fam\) est génératrice \(\Longleftrightarrow\) \(\Fam\) est une base.
}
En effet, intuitivement on comprends que si une famille possède \textbf{le bon nombre de vecteurs}, alors il suffit qu'elle soit libre (ou génératrice) pour que ce soit une base.\<

Soient \(F, G\) deux sous-espaces de \(E\). La dimension permet aussi de prouver des \textbf{égalités} d'espaces vectoriels, gràce à la propriété:
\[
   \begin{rcases}
      F \subset G \\
      \text{dim}(F) = \text{dim}(G)
   \end{rcases}\implies F = G
\]
Aussi, on peut alors \textbf{caractériser} le fait que \(F\) et \(G\) soient \textbf{supplémentaires} dans \(E\) par:
\[
   F \oplus G = E \Longleftrightarrow \begin{cases}
      F \cap G = \{0_E\}\\
      \text{dim}(F) + \text{dim}(G) = \text{dim}(E)
   \end{cases}
\]
Enfin on peut calculer \textbf{la dimension d'une somme}\footnote[1]{Le cas d'égalité est intéressant car il intervient pour des sous-espaces d'intersection nulle, et alors la dimension de la somme est simplement égale à la somme des dimensions.} avec la \textbf{formule de Grassmann}:
\customBox{width=8cm}{
   \(
      \text{dim}(F + G) = \text{dim}(F) + \text{dim}(G) - \text{dim}(F \cap G)
   \)
}


\subsection*{\subsecstyle{Rang {:}}}
\addcontentsline{toc}{subsection}{Rang}

Soit \(E\) un espace vectoriel de dimension finie et \(\Fam = (e_i)_{i \in \N}\) une famille de vecteurs de cet espace, alors on appelle \textbf{rang} de \(\Fam\) l'entier:
\customBox{width=5cm}{
   \(\text{rg}(\Fam) = \text{dim}(\text{Vect}(\Fam))\)
}

Informellement, le rang est simplement \textbf{la dimension de l'espace engendré par la famille}.\<

On étends cette définition à celle du \textbf{rang d'une application linéaire}, qu'on note \(\text{rg}(f)\), qui correspond à \textbf{la dimension de son image}, et on définit ainsi une application linéaire \textbf{de rang fini}.

\subsection*{\subsecstyle{Applications linéaires en dimension finie{:}}}
\addcontentsline{toc}{subsection}{Applications linéaires en dimension finie}

Soit \(f \in \mathcal{L}(E, F)\) une telle application, alors on peut montrer le \textbf{théorème du rang}:
\customBox{width=7cm}{
   \(
      \text{dim}(E) = \text{rg}(f) + \text{dim}(\text{Ker}(f))
   \)
}

Ce thèorème permet alors de caractériser l'injectivité et la surjectivité d'une application linéaire par:
\begin{itemize}
   \item \(f\) est injective si et seulement si \(\text{rg}(f) = \text{dim}(E)\)
   \item \(f\) est surjective si et seulement si \(\text{rg}(f) = \text{dim}(F)\)
\end{itemize}

En particulier, si \(E\) et \(F\) sont deux espaces vectoriels \textbf{de même dimension} alors:
\customBox{width=10cm}{
   \(f\) est injective \(\Longleftrightarrow\) \(f\) est surjective \(\Longleftrightarrow\) \(f\) est bijective
}
On caractérise ainsi \textbf{les isomorphismes en dimension finie}\footnote[2]{Toutes ces propriétés n'étant que de simples applications du théorème du rang.}.

\chapter*{\chapterstyle{II --- Espace des matrices}} % Fini 99%
\addcontentsline{toc}{section}{Espace des matrices}

On appelle \textbf{matrice} à \(n\) lignes et \(p\) colonnes à coefficients dans \(\K\) toute application \(M : \inticc{1}{n} \times \inticc{1}{p} \longrightarrow \K\), en d'autres termes \textbf{tout famille de scalaires indexée par }\(\inticc{1}{n} \times \inticc{1}{p}\).\<

Informellement, il s'agit d'une généralisation du concept de suite sous forme de suite à \textbf{deux indices}, qu'on peut alors voir comme un tableau de nombres tel qu'en chaque position \((i, j)\), on ait un scalaire \(a_{ij} \in \K\).\< 

A l'instar des suites, on note \(M = (a_{ij})_{\substack{1 \leq i \leq n\\1 \leq j \leq p}}\) pour faire référence au terme général de \(M\) en position \((i,j)\).

\subsection*{\subsecstyle{Définition et propriétés{:}}}
\addcontentsline{toc}{subsection}{Définition}

On appelle \textbf{espace des matrices} à \(n\) lignes et \(p\) colonnes à coefficient dans \(\K\), et on note \(\mathcal{M}_{n, p}(\K)\), l'espace de toutes les matrices comme définies ci-dessus.\<

On définit \textbf{la somme} de deux matrices comme la matrice dont le terme en position \((i, j)\) est la sommes des termes des deux matrices en cette position, informellement, on additionne simplement \textbf{composante par composante}.\+
On définit \textbf{la multiplication} d'une matrice par un scalaire comme la matrice dont tout les termes sont multipliés par ce scalaire.\<

Pour ces deux lois l'ensemble \((\mathcal{M}_{n, p}(\K), +, \cdot)\) est alors muni d'une structure \textbf{d'espace vectoriel} sur \(\K\) (de dimension \(n \times p\)) ayant pour vecteur nul \(0_E\) la matrice dont tout les coefficients sont nuls.

\subsection*{\subsecstyle{Produit matriciel{:}}}
\addcontentsline{toc}{subsection}{Produit matriciel}

Soit \(A = (a_{i, j}) \in \mathcal{M}_{n, m}(\K)\) et \(B = (b_{k, j}) \in \mathcal{M}_{m, p}(\K)\).\+
Alors on définit\footnote[1]{Il faut que le nombre de colonnes de la première soit égal au nombre de lignes de la seconde pour que les matrices soient dites \textbf{compatibles}.}  la matrice \(C := AB = (c_{i,j})\) comme étant la matrice telle que:
\customBox{width=3.5cm}{
   \[
      c_{i, j} = \sum_{k=1}^{n}a_{i, k}b_{k, j}
   \]
}

En particulier, on appelle ce produit \textbf{un produit ligne par colonne} qui se comprends visuellement par:\+
\begin{center}
   \begin{tikzpicture}[    
      every left delimiter/.style={xshift=.55em},
      every right delimiter/.style={xshift=-.55em}
    ]
      \matrix[
        matrix of math nodes, ampersand replacement=\&,
        left delimiter=(, right delimiter=),  outer sep = 0pt,inner sep=4.5pt
      ](A){
        a_{1, 1} \& a_{1, 2}\\ 
        a_{2, 1} \& a_{2, 2}\\ 
        a_{3, 1} \& a_{3, 2}\\  
      };
      \draw[color=BrightBlue1, line width=0.5mm] (A-3-1.south west) rectangle (A-3-2.north east);
      \draw node at (1.5, 0) {\(\times\)};
  
      \matrix[
        matrix of math nodes, ampersand replacement=\&,
        left delimiter=(, right delimiter=),  outer sep = 0pt,inner sep=4.5pt
      ] (B) at (3.4, 0) {
        b_{1, 1} \& b_{1, 2} \& b_{1, 3}\\ 
        b_{2, 1} \& b_{2, 2} \& b_{2, 3}\\ 
      };
      \draw[color=BrightBlue1, line width=0.5mm] (B-1-1.north west) rectangle (B-2-1.south east);
  
      \draw node at (5.25, 0) {\(=\)};
  
      \matrix[
        matrix of math nodes, ampersand replacement=\&,
        left delimiter=(, right delimiter=),  outer sep = 0pt,inner sep=4.5pt
      ] (C) at (7.1, 0) {
        c_{1, 1} \& c_{1, 2} \& c_{1, 3}\\ 
        c_{2, 1} \& c_{2, 2} \& c_{2, 3}\\ 
        c_{3, 1} \& c_{3, 2} \& c_{3, 3}\\
      };
      \draw[color=BrightRed1, line width=0.5mm] (C-3-1.north west) rectangle (C-3-1.south east);
    \end{tikzpicture}
\end{center}
Le coefficient à \color{BrightRed1} la troisième ligne, première colonne \color{black} est obtenu en multipliant \color{BrightBlue1} la troisième ligne par la première colonne.\color{black}\<

On appelle \textbf{matrice identité} la matrice neutre pour cette opération, c'est la matrice composée uniquement de \(1\) sur toute la diagonale principale, et de \(0\) partout ailleurs.\<

En considérant cette définition, il est évident que le produit matrice n'est en général \textbf{pas commutatif}. Par ailleurs \textbf{le produit de deux matrices non nulles peut être nul}.
\pagebreak

Sous réserve de compatibilités des tailles matrices, on a les propriétés suivantes:
\begin{multicols}{2}
   \begin{itemize}
      \item Associativité
      \item La matrice identité est neutre
      \item Distributivité
      \item La matrice nulle est absorbante
   \end{itemize}
\end{multicols}   

Considérons maintenant des matrices \textbf{carrées} de même taille, le produit est alors toujours défini et on peut définir une \textbf{puissance positive de matrice}, qu'on note \(A^n\) qui vérifie alors:

\begin{multicols}{2}
   \begin{itemize}
      \item \((A^p)^q = A^{pq}\)
      \item \(A^pA^q = A^{p +q}\)
   \end{itemize}
\end{multicols}    

\subsection*{\subsecstyle{Transposition{:}}}
\addcontentsline{toc}{subsection}{Transposition}

Soit \(M = (x_{i,j})\in \mathcal{M}_{n, m}(\K)\), on définit maintenant \textbf{l'opération de transposition} d'une matrice notée \(M^\mathsf{T}\), c'est \textbf{une application linéaire involutive} définie par:
\customBox{width=3cm}{
   \(
      M^\mathsf{T} = (x_{j,i})
   \)
}

Intuitivement, cette application transforme \textbf{chaque ligne en colonne et invérsement}. Par exemple:
\[
   \begin{pmatrix}
      a & b & c\\
      d & e & f
   \end{pmatrix}^\top = 
   \begin{pmatrix}
      a & d\\
      b & e\\
      c & f
   \end{pmatrix}
\]
Il faut aussi noter son comportement par rapport au produit matriciel de deux matrices \(A, B\), précisément on a:
\customBox{width=3.5cm}{
   \(
      (AB)^\top = B^\top A^\top
   \)
}


\subsection*{\subsecstyle{Trace{:}}}
\addcontentsline{toc}{subsection}{Trace}

On définit aussi une autre application linéaire appellée \textbf{trace d'une matrice}, et qui est définie comme \textbf{la somme des éléments diagonaux}. Formellement:
\[
   \text{tr}(A) := \sum_{i=1}^{n} a_{i, i}   
\]
Elle est donc linéaire mais on a aussi:
\customBox{width=3.5cm}{
   \(
      \text{tr}(AB) = \text{tr}(BA)  
   \)
}


\subsection*{\subsecstyle{Matrice d'une famille de vecteurs{:}}}
\addcontentsline{toc}{subsection}{Matrice d'une famille de vecteurs}

Soit \(\mathscr{F} = (e_i)_{i\in\N}\), alors on peut définir \textbf{la matrice de la famille} dans une base \(\mathscr{B}\) par:
\customBox{width=7cm}{
   \(
      \text{Mat}_{\mathscr{B}}(\mathscr{F}) = ([e_1]_{\mathscr{B}}, [e_2]_{\mathscr{B}}, \ldots, [e_n]_{\mathscr{B}})  
   \)
}

\begin{center}
   \textit{
      La matrice d'une famille est donc constituée des coordonées des vecteurs\+
      dans la base donnée (en colonnes).
   }
\end{center}
\pagebreak

\subsection*{\subsecstyle{Matrice d'une application linéaire{:}}}
\addcontentsline{toc}{subsection}{Matrice d'une application linéaire}

Soit \(E, F\) deux espaces vectoriels (de bases \(\mathscr{B} = (e_i)_{i \in \N}\) et \(\mathscr{C}=(f_i)_{i \in \N}\)) et \(f \in \mathcal{L}(E, F)\).\<

Alors on peut associer à l'application \(f\) \textbf{une unique matrice} dans les bases \(\mathscr{B, C}\), qu'on note alors Mat\(_{\mathscr{B, C}}(f)\) qu'on construit comme suit:
\customBox{width=9cm}{
   \(
      \text{Mat}_{\mathscr{B, C}}(f) = ([f(e_1)]_{\mathscr{C}}, [f(e_2)]_{\mathscr{C}}, \ldots, [f(e_n)]_{\mathscr{C}})
   \)
}
\begin{center}
   \textit{
      La matrice d'une application linéaire est donc constituée des coordonées dans la base d'arrivée de l'image des vecteurs de la base de départ (en colonnes).
   }
\end{center}

\underline{Exemple:} Considérons un endomorphisme de \(\R^3\) et sa base canonique notée \((e_1, e_2, e_3)\), telle que \(f(x, y, z) = (2x+1, 3x, z-1)\). Alors on calcule l'image des vecteurs de la base de départ, ie:
\begin{multicols}{3}
\begin{itemize}
   \item \(f(1, 0, 0) = (3, 3, -1)\)
   \item \(f(0, 1, 0) = (0, 0, -1)\)
   \item \(f(0, 0, 1) = (0, 0, -2)\)
\end{itemize}
\end{multicols}
Puis on calcule les coordonées de ces vecteurs dans la base d'arrivée, et on les range en colonne dans une matrice et on obtient:
\begin{center}
   \begin{tikzpicture}[      
      every left delimiter/.style={xshift=.55em},
      every right delimiter/.style={xshift=-.55em}
    ]
    \draw node[color = BrightBlue1] at (-1.05, 1.75) {$f(e_1)$};
    \draw[-latex, dashed, thick, color = BrightBlue1!50] (-1.05, 1.5) -- (-1.05, 1.1);
  
    \draw node[color = BrightBlue1] at (0, 1.75) {$f(e_2)$};
    \draw[-latex, dashed, thick, color = BrightBlue1!50] (0, 1.5) -- (0, 1.1);
  
    \draw node[color = BrightBlue1] at (1.05, 1.75) {$f(e_3)$};
    \draw[-latex, dashed, thick, color = BrightBlue1!50] (1.05, 1.5) -- (1.05, 1.1);
  
    \draw node[color = BrightBlue1] at (2.3, 0.85) {$f_1$};
    \draw[-latex, dashed, thick, color = BrightBlue1!50] (1.7, 0.85) -- (2.1, 0.85);
  
    \draw node[color = BrightBlue1] at (2.3, 0) {$f_2$};
    \draw[-latex, dashed, thick, color = BrightBlue1!50] (1.7, 0) -- (2.1, 0);
  
    \draw node[color = BrightBlue1] at (2.3, -0.85) {$f_3$};
    \draw[-latex, dashed, thick, color = BrightBlue1!50] (1.7, -0.85) -- (2.1, -0.85);
  
    \matrix[
      matrix of math nodes, ampersand replacement=\&,
      left delimiter=(, right delimiter=), inner ysep=4pt, column sep=0.8em, row sep = 0.8em,
      nodes={
        inner sep=0.5em,outer sep=0,text width=1.2em,align=center
      }
    ]{
      3 \& 0 \& 0\\ 
      3 \& 0 \& 0\\ 
      -1 \& -1 \& -2\\
    };
    \end{tikzpicture}
\end{center}


Cette construction implique qu'il suffit donc, pour une base donnée, de \textbf{calculer les images des vecteurs de cette base}, puis leur coordonées pour décrire la transformation.\<

Si \(E\) et \(F\) sont de dimensions respectives \(n\) et \(p\), on peut alors construire \textbf{un isomorphisme fondamental}:
\customBox{width=7cm}{
   \(
      \text{Mat}_{\mathscr{B, C}} : \mathcal{L}(E, F) \xlongrightarrow{\sim} \mathcal{M}_{n, p}(\K) 
   \)
}

\begin{center}
   \textit{On peut alors considèrer, que du point de vue de l'algèbre linéaire, l'espace des application linéaires et l'espace des matrices sont identiques\footnote[1]{En dimension finie}.}
\end{center}

Dans la suite, pour plus de lisibilité, on considèrera un endomorphisme \(f\) de \(\R^2\) de matrice \(M\) dans la base canonique \(\mathscr{C}\).
Alors, appliquer une transformation linéaire à un vecteur \(u \in E\) revient à \textbf{multiplier ce vecteur par la matrice associée}\footnote[2]{Il suffit de partir du produit matriciel et d'utiliser la définition de \(M\) comme étant la matrice de \(f\) et les règles de calculs sur les vecteurs colonnes.}, ie on a:
\customBox{width=4cm}{
   \(
      [f(u)]_{\mathscr{C}} = M \times [u]_{\mathscr{C}}
   \)
}

\underline{Exemple:} Soit \(f\) de matrice 
\(M = \begin{pmatrix}
   1 & 2\\
   3 & 4\\
\end{pmatrix}\) 
et \(u = (5, 3)\), calculer les coordonées de \(f(u)\) revient à calculer:
\[
   \begin{pmatrix}
      1 & 2\\
      3 & 4\\
   \end{pmatrix}
   \begin{pmatrix}
      5 \\
      3 \\
   \end{pmatrix} =
   \begin{pmatrix}
      11 \\
      27 \\
   \end{pmatrix} 
\]
Enfin l'isomorphisme entre ces deux espaces nous permet de définir \textbf{le noyau et l'image d'une matrice}, comme simplement étant le noyau ou l'image de l'application linéaire associée à cette matrice.

\subsection*{\subsecstyle{Matrices inversibles{:}}}
\addcontentsline{toc}{subsection}{Matrices inversibles}

Soit \(M \in \mathcal{M}_n(\K)\) une matrice carrée, alors on dit que \(M\) est \textbf{inversible} si il existe un inverse pour la loi de multiplication des matrices, ie si il existe une matrice \(A^{-1}\) telle que:
\customBox{width=3.5cm}{
   \(
      AA^{-1}=\text{Id}_n
   \)
}

En particulier, si on considère l'application linéaire associée à \(M\), alors:
\begin{center}
   \(M\) inversible \(\Longleftrightarrow\) \(f\) est inversible
\end{center}

On appelle l'ensemble des matrices carrées inversibles de taille \(n\) \textbf{le groupe linéaire} d'ordre \(n\), qu'on note GL\(_n(\K)\), c'est un groupe pour la multiplication matricielle, ie on a:
\[
   A, B \in \text{GL}_n(\K) \Longleftrightarrow AB \in \text{GL}_n(\K)
\]
Et en particulier, on montre\footnote[1]{Intuitivement, c'est équivalent à l'inverse d'une composée car les matrices ``sont'' des applications au sens de l'algèbre linéaire.} que \((AB)^{-1} = B^{-1}A^{-1}\). Dans le cas d'une matrice inversible, on peut alors étendre notre définition de puissance d'une matrice au cas d'entiers négatifs.\<

Intuitivement:
\begin{center}
   \textit{Le groupe \(\text{GL}_n(\K)\) est donc le groupe \textbf{des automorphismes} de \(E\), c'est à dire le groupe \textbf{des transformations} (linéaires) de \(E\) qui sont \textbf{réversibles}.}
\end{center}
Visuellement, on comprends directement l'idée d'automorphisme d'un espace, son lien avec la bijectivité, la matrice inverse est alors la représentation algèbrique de la transformation inverse (ici avec une rotation de \(\R^2\) de \(45\) degrés):
\begin{center}
   \begin{tikzpicture}
      \draw[step=0.45cm, color=black!20, dashed] (-1.75,-1.75) grid (1.75,1.75);
      \draw[-latex, color=BrightBlue1, thick] (0,0) -- (0,1.5) node[below left] {$y$};
      \draw[-latex, color=BrightBlue1, thick] (0,0) -- (1.5,0) node[below left] {$x$};
      \draw[-latex, color=BrightRed1, thick] (0,0) -- (1,1.2) node[below right] {$u$};
  
      \draw[-latex] (2,0) -- (3,0);
      \draw[] (2,-0.05) -- (2,0.05);
      \draw[] node[] at (2.5,0.3){$\varphi$};
  
      \draw[step=0.45cm, color=black!20, dashed, shift={(5.5 cm, 0)}, rotate=45] (-1.75,-1.75) grid (1.75,1.75);
      \draw[-latex, color=BrightBlue1, thick, shift={(5.5 cm, 0)}, rotate=45] (0,0) -- (0,1.5) node[below left] {$\varphi(y)$};
      \draw[-latex, color=BrightBlue1, thick, shift={(5.5 cm, 0)}, rotate=45] (0,0) -- (1.5,0) node[below right] {$\varphi(x)$};
      \draw[-latex, color=BrightRed1, thick, shift={(5.5 cm, 0)}, rotate=45] (0,0) -- (1,1.2) node[below right] {$\varphi(u)$};
  
      \draw[-latex] (8,0) -- (9,0);
      \draw[] (8,-0.05) -- (8,0.05);
      \draw[] node[] at (8.65,0.375){$\varphi^{-1}$};
  
      \draw[step=0.45cm, color=black!20, dashed, shift={(11 cm, 0)}] (-1.75,-1.75) grid (1.75,1.75);
      \draw[-latex, color=BrightBlue1, thick, shift={(11 cm, 0)}] (0,0) -- (0,1.5) node[below left] {$y$};
      \draw[-latex, color=BrightBlue1, thick, shift={(11 cm, 0)}] (0,0) -- (1.5,0) node[below left] {$x$};
      \draw[-latex, color=BrightRed1, thick, shift={(11 cm, 0)}] (0,0) -- (1,1.2) node[below right] {$u$};
    \end{tikzpicture}   
\end{center}
Enfin, on peut montrer que \textbf{la transposition} est compatible avec l'inverse, ie on a:
\customBox{width=4cm}{
   \(
      (M^{-1})^{\top} = (M^{\top})^{-1} 
   \)
}  
Montrer qu'une matrice est inversible consiste simplement à calculer \textbf{le rang de la matrice}, comme expliqué dans la partie suivante.\+ 
Par ailleurs, calculer l'inverse peut se réaliser de plusieurs manières, la méthode la plus générale consiste à résoudre l'équation d'inconnue \(X\) de la forme:
\[
   AX = Y   
\]

On définira aussi une application fondamentale nommée \textbf{déterminant} qui sera définie dans le prochain chapitre et qui caractérise exactement toutes les matrices inversibles et donne une autre manière de calculer l'inverse.
\pagebreak

\subsection*{\subsecstyle{Rang d'une matrice{:}}}
\addcontentsline{toc}{subsection}{Rang d'une matrice}

On a défini précédemment le concept de rang d'un famille et de rang d'une application linéaire, on peut généraliser cette définition en celle de \textbf{rang d'une matrice}, en effet on a simplement:
\begin{center}
   \textit{Le rang d'une matrice est la dimension de l'espace de ses colonnes.}
\end{center}

Soit \(M \in \mathcal{M}_n(\K)\), dont on note les colonnes \(\Fam = (C_1, C_2, \ldots, C_n)\), alors si le rang de la matrice est \(n\), on peut en déduire plusieurs propriétés:
\begin{multicols}{2}
   \begin{itemize}
      \item La matrice est inversible.
      \item La famille \(\Fam\) est une base de \(\K^n\).
   \end{itemize}
\end{multicols}
On remarque donc que le rang nous donne \textbf{énormément d'informations} sur la matrice, la famille des colonnes, et l'application linéaire associée.

\subsection*{\subsecstyle{Matrice de passage{:}}}
\addcontentsline{toc}{subsection}{Matrice de passage}

On considère un espace vectoriel \(E\) et deux bases \(\Fam=(e_1, \ldots, e_n), \Fam'=({e}_1', \ldots, {e}_n')\) de cet espace.\+
On appelle \textbf{matrice de passage} de \(\Fam\) à \(\Fam'\) la matrice \(P\) qu'on construit comme suit:
\customBox{width=7cm}{
   \(
      \text{Pass}(\Fam, \Fam') = ([e_1']_{\Fam}, [e_2']_{\Fam}, \ldots, [e_n']_{\Fam}) 
   \)
} 
\begin{center}
   \textit{
      La matrice de passage est donc constituée des coordonées dans l'ancienne base des vecteurs de la nouvelle base(en colonnes).
   }
\end{center}

\underline{Exemple:} Considèrons deux bases de \(\R^2\) telles que \(\Fam = [(1, 2), (3, 4)]\) et \(\Fam' = [(5, 6), (7, 8)]\), alors on a:
\[
   [5, 6]_\Fam = \begin{pmatrix} -1 \\ 2\end{pmatrix}
   \quad\quad\quad\quad
   [7, 8]_\Fam = \begin{pmatrix} 3 \\ -2\end{pmatrix}
\]
On range ensuite ces coordonées en colonnes dans la matrice, et on obtient:
\begin{center}
   \vspace{-5pt}   
   \begin{tikzpicture}[      
      every left delimiter/.style={xshift=.55em},
      every right delimiter/.style={xshift=-.55em},
   ]
   \draw node[] at (-2.35, 0) {$\text{Pass}(\Fam, \Fam') = $};

   \draw node[color = BrightBlue1] at (-0.5, 1.5) {$e_1'$};
   \draw[-latex, dashed, thick, color = BrightBlue1!50] (-0.5, 1.25) -- (-0.5, 0.85);

   \draw node[color = BrightBlue1] at (0.5, 1.5) {$e_2'$};
   \draw[-latex, dashed, thick, color = BrightBlue1!50] (0.5, 1.25) -- (0.5, 0.85);

   \draw node[color = BrightBlue1] at (1.75, 0.440) {$e_1$};
   \draw[-latex, dashed, thick, color = BrightBlue1!50] (1.1, 0.465) -- (1.5, 0.465);

   \draw node[color = BrightBlue1] at (1.75, -0.490) {$e_2$};
   \draw[-latex, dashed, thick, color = BrightBlue1!50] (1.1, -0.465) -- (1.5, -0.465);

   \matrix[
      matrix of math nodes, ampersand replacement=\&,
      left delimiter=(, right delimiter=), inner ysep=2pt, column sep=0.8em, row sep = 0.8em,
      nodes={
         inner sep=0.5em,outer sep=0,text width=1em,align=center
      }
   ]{
      -1 \& 3 \\ 
      2 \& -2 \\ 
   };
   \end{tikzpicture}
\end{center}

Une matrice de passage est alors \textbf{nécessairement inversible}\footnote[1]{Car son rang est égal à sa dimension par construction.}, et si on note \(U = [u]_{\Fam}\) et \(U' = [u]_{\Fam'}\)on a alors le théorème fondamental suivant:
\customBox{width=3cm}{
   \(
      U' = P^{-1}U
   \)
} 

\begin{center}
   \textit{La matrice de passage nous permet alors de représenter un même vecteur\+ 
   dans une base différente.}
\end{center}
Considérons maintenant un endomorphisme \(f\) de matrices \(M = \text{Mat}(\Fam, f), M' = \text{Mat}(\Fam', f)\) dans les deux bases respectivement, alors on a:
\customBox{width=3.5cm}{
   \(
      M' = P^{-1}MP
   \)
}
\begin{center}
   \textit{La matrice de passage nous permet alors de représenter un même endomorphisme\+ 
   dans une base différente.}
\end{center}
\pagebreak

\subsection*{\subsecstyle{Matrices semblables{:}}}
\addcontentsline{toc}{subsection}{Matrices semblables}

On dira alors que deux matrices \(A, B \in \mathcal{M}_n(\K)\) sont \textbf{semblables} si il existe une matrice de passage telle que la relation ci-dessus soit vérifiée.
Ce qui signifie exactement que:
\customBox{width=14cm}{
   \textit{Deux matrices semblables représentent \textbf{la même transformation géométrique} \+
   dans des bases différentes.}
}

En particulier on appelle alors \textbf{invariants de similitude} les propriétés qui \textbf{ne dépendent pas du choix de la base}, on peut alors montrer que:
\customBox{width=13cm}{
   Le rang et la trace d'une matrice sont des invariants de similitude.
}
   
Un cas remarquable de matrices semblables est celui de la transposée, en effet dans un corps, une matrice est sa transposée sont semblables.

\subsection*{\subsecstyle{Matrices remarquables{:}}}
\addcontentsline{toc}{subsection}{Matrices remarquables}

Dans ce chapitre nous allons présenter brièvement différentes matrices remarquables:
\begin{figure}[h]
   \color{DarkBlue1}
   \centering
   \begin{subfigure}{.3\textwidth}
      \centering
      \begin{tikzpicture}[line cap=round]
         \matrix[
            matrix of math nodes, ampersand replacement=\&,
            left delimiter=(, right delimiter=)
         ]{
            a \& 0 \& 0\\ 
            0 \& d \& 0\\ 
            0 \& 0 \& f\\
         };
      \end{tikzpicture}
      \caption*{\color{DarkBlue1}
      Matrice diagonale}
   \end{subfigure}\quad
   \begin{subfigure}{.3\textwidth}
      \centering
      \begin{tikzpicture}[line cap=round]
         \matrix[
            matrix of math nodes, ampersand replacement=\&,
            left delimiter=(, right delimiter=)
         ]{
            a \& b \& c\\ 
            0 \& d \& e\\ 
            0 \& 0 \& f\\
         };
      \end{tikzpicture}
      \caption*{\color{DarkBlue1}
      Matrice triangulaire supérieure}
   \end{subfigure}\quad
   \begin{subfigure}{.3\textwidth}
      \centering
      \begin{tikzpicture}[line cap=round]
         \matrix[
            matrix of math nodes, ampersand replacement=\&,
            left delimiter=(, right delimiter=)
         ]{
            a \& 0 \& 0\\ 
            b \& d \& 0\\ 
            c \& e \& f\\
         };
      \end{tikzpicture}
      \caption*{\color{DarkBlue1}
      Matrice triangulaire inférieure}
   \end{subfigure}
\end{figure}

Puis on a deux types de matrices qui sont liées à la symétrie des coefficients et à \textbf{l'opération de transposition}\footnote[1]{En effet une matrice symétrique est égale à sa transposée et une matrice antisymétrique à \textbf{l'opposée} de sa transposée}:
\begin{figure}[h]
   \color{DarkBlue1}
   \centering
   \begin{subfigure}{.3\textwidth}
      \centering
      \begin{tikzpicture}[line cap=round]
         \matrix[
            matrix of math nodes, ampersand replacement=\&,
            left delimiter=(, right delimiter=)
         ]{
            a \& d \& f\\ 
            d \& b \& e\\ 
            f \& e \& c\\
         };
      \end{tikzpicture}
      \caption*{\color{DarkBlue1}Matrice symétrique}
   \end{subfigure}\quad
   \begin{subfigure}{.3\textwidth}
      \centering
      \begin{tikzpicture}[line cap=round]
         \matrix[
            matrix of math nodes, ampersand replacement=\&,
            left delimiter=(, right delimiter=)
         ]{
            0 \& -a \& -b\\ 
            a \& 0 \& -c\\ 
            b \& c \& 0\\
         };
      \end{tikzpicture}
      \caption*{\color{DarkBlue1}Matrice antisymétrique}
   \end{subfigure}
\end{figure}

On a aussi \textbf{les matrices élémentaires} qui sont obtenus par \textbf{opérations élémentaires}\footnote[2]{Voir la partie suivante.} sur la matrice identité:
\begin{figure}[h]
   \color{DarkBlue1}
   \centering
   \begin{subfigure}{.3\textwidth}
      \centering
      \begin{tikzpicture}[line cap=round]
         \matrix[
            matrix of math nodes, ampersand replacement=\&,
            left delimiter=(, right delimiter=)
         ]{
            1 \& 0 \& 0\\ 
            0 \& 1 \& 0\\ 
            0 \& 0 \& \lambda\\
         };
      \end{tikzpicture}
      \caption*{\color{DarkBlue1}Matrice de dilatation}
   \end{subfigure}\quad
   \begin{subfigure}{.3\textwidth}
      \centering
      \begin{tikzpicture}[line cap=round]
         \matrix[
            matrix of math nodes, ampersand replacement=\&,
            left delimiter=(, right delimiter=)
         ]{
            1 \& 0 \& 0\\ 
            0 \& 1 \& 0\\ 
            0 \& \lambda\& 1\\
         };
      \end{tikzpicture}
      \caption*{\color{DarkBlue1}Matrice de transvection}
   \end{subfigure}\quad
   \begin{subfigure}{.3\textwidth}
      \centering
      \begin{tikzpicture}[line cap=round]
         \matrix[
            matrix of math nodes, ampersand replacement=\&,
            left delimiter=(, right delimiter=)
         ]{
            1 \& 0 \& 0\\ 
            0 \& 0 \& 1\\ 
            0 \& 1\& 0\\
         };
      \end{tikzpicture}
      \caption*{\color{DarkBlue1}Matrice de permutation}
   \end{subfigure}
\end{figure}

Enfin, on a \textbf{les matrices orthogonales} qui sont des matrices telles que leur inverse soit \textbf{leur transposée}, ce sont par exemple les matrices de rotation:
\begin{figure}[H]
   \color{DarkBlue1}
   \centering   
   \begin{subfigure}{.3\textwidth}
      \centering
      \begin{tikzpicture}[line cap=round]
         \matrix[
            matrix of math nodes, ampersand replacement=\&,
            left delimiter=(, right delimiter=)
         ]{
            \cos(\theta) \& -\sin(\theta)\\ 
            \sin(\theta) \& \cos(\theta) \\
         };
      \end{tikzpicture}
      \caption*{\color{DarkBlue1}Matrice de rotation}
   \end{subfigure}
\end{figure}

\pagebreak
\subsection*{\subsecstyle{Echelonnements \& Calculs{:}}}
\addcontentsline{toc}{subsection}{Echelonnements et calculs}

La théorie permettant de lier matrices, applications linéaires et familles de vecteurs, l'étude des espaces vectoriels engendrés par les colonnes ou les lignes d'une matrice revêt alors une importantce capitale qui est l'objet de cette partie.\<

On appelle \textbf{opérations élémentaires} sur une matrice \(M\), l'une de ces 3 opérations:
\begin{itemize}
   \item Multiplier une colonne par un scalaire
   \item Ajouter une colonnes à une autre
   \item Echanger deux colonnes
\end{itemize}
Il est alors possible de montrer que ces trois opérations préservent \textbf{le sous-espace engendré}\footnote[1]{L'échange ne change évidement rien, et le sous espace engendré ne change pas pas multiplication par un scalaire ou ajout d'un vecteur venant de la famille en question.} par les colonnes de la matrice, donc en particulier \textbf{le rang, le noyau et l'image} sont préservés.\<

Enfin, on appellera \textbf{matrice équivalente} à \(M\) et on notera \(M' \sim M\) la matrice \(M\) à laquelle on a appliqué une ou plusieurs opérations élémentaires.

On dira qu'une matrice \(M\) est \textbf{échelonée} si la matrice obtenue aprés application d'opérations élémentaires est de la forme générale:
\begin{center}
   \begin{tikzpicture}[    
      every left delimiter/.style={xshift=.55em},
      every right delimiter/.style={xshift=-.55em}
   ]
      \matrix[
         matrix of math nodes, ampersand replacement=\&,
         left delimiter=(, right delimiter=),  outer sep = 0pt,inner sep=4.5pt
      ] (A) {
         {\alpha_1} \& {0} \& {0} \& {0}\& {0}\\ 
         {*} \& {\alpha_2} \& {0} \& {0}\& {0}\\ 
         {*} \& {*} \& {0} \& {0}\& {0}\\ 
         {*} \& {*} \& {\alpha_3} \& {0}\& {0}\\
      };
      \draw[color=BrightBlue1, line width=0.5mm] (A-1-1.north west) -- (A-1-1.north east) -- (A-1-1.south east) -- (-0.1, 0.57)-- (-0.1, -0.64) -- (A-4-3.north east)-- (A-4-3.south east);

   \end{tikzpicture}   
\end{center}
En d'autres termes, on \textbf{creuse} la matrice pour faire apparaître des zéros sur la partie triangulaire supérieure.\+
Il existe aussi une variante appellé \textbf{échelonnement avec mémorisation}, pour cette variante, on nomme chaque vecteur-colonne de la matrice et on reporte toutes les transformations réalisées sur ces vecteurs. \<

Présentons maintenant les différentes informations que nous pouvons tirer de ces échellonements:

\begin{enumerate}
   \item On peut trés facilement trouver le rang d'une matrice, en effet aprés échelonnement, c'est simplement \textbf{le nombre de colonnes non-nulles} de la matrice. En particulier, il est donc trés facile de montrer qu'une matrice est inversible ou qu'une famille est une base.
   \item On peut trouver une base d'un sous-espace engendré par une famille, pour cela on échelonne la matrice de cette famille, l'ensemble des colonnes non nulles consitue une base du sous-espace initial.
   \item On peut trouver un supplémentaire d'une famille, on échelonne la matrice de la famille, et on complète la matrice par des vecteurs de la base canonique\footnote[2]{Ou alors de telle sorte que la matrice reste échelonnée.}, ces vecteurs formeront alors un supplémentaire.
   \item On peut trouver les équations cartésiennes d'un sous-espace engendré par une famille, pour cela on augmente la matrice de la famille par une matrice d'indéterminées, et on échelonne. Alors la dernière colonne fournira les equations cartésiennes du sous-espace.
   \item On peut trouver le \textbf{noyau, l'image} d'une matrice, gràce à la variante \textbf{avec mémorisation}, alors les colonnes nulles aprés échelonnement permettent d'en déduire des vecteurs envoyés du noyau, et les colonnes non-nulles constituent une base de l'image (voir l'exemple ci-aprés).
\end{enumerate}
Il peut être utile de noter que fondamentalement, on peut à la fois échelonner \textbf{sur les colonnes} ou sur \textbf{les lignes}. En effet, le rang est invariant par transposition, on aurait tout aussi bien pu développer des opération élémentaires sur les lignes, et les résultats seraient équivalents.

\pagebreak

Soit \(M = (C_1, C_2, C_3) = \begin{pmatrix} 1 & 2 & 1\\ 3 & 4 & 2\\ 5 & 6 & 3\end{pmatrix}\), nous allons présenter quelques exemples:
\begin{itemize}
   \item Cherchons le rang de la matrice, on a:
   \[
      M = \begin{pmatrix} 1 & 2 & 1\\ 3 & 4 & 2\\ 5 & 6 & 3\end{pmatrix} \sim  \begin{pmatrix} 1 & 0 & 0\\ 3 & 2 & 2\\ 5 & 4 & 4\end{pmatrix} \sim  \begin{pmatrix} 1 & 0 & 0\\ 3 & 2 & 0\\ 5 & 4 & 0\end{pmatrix} \sim  \begin{pmatrix} 1 & 0\\ 3 & 2\\ 5 & 4\end{pmatrix}
   \]
   La matrice est échelonnée sur \(2\) colonnes, donc son rang est \(2\). En particulier, \((C_1, C_2, C_3)\) n'est pas une base de \(\R^3\), la matrice (donc l'endomorphisme associé) n'est pas inversible.
   \item Cherchons un sous-espace supplémentaire à \((C_1, C_2, C_3)\) en reprenant l'échelonnement ci-dessus, on a simplement à rajouter des colonnes à \(M\) de sorte qu'elle soit échelonnée sur 3 colonnes, par exemple:
   \[
      M' = \begin{pmatrix} 1 & 0 & 0\\ 3 & 2 & 0\\ 5 & 4 & 1\end{pmatrix}
   \]
   Et donc le sous-espace engendré par le vecteur \((0, 0, 1)\) est bien un supplémentaire de l'espace des colonnes.
   \item Cherchons des équations cartésiennes du sous-espace des colonnes, on a:
   \[
      M = \begin{pmatrix} 1 & 2 & 1 & x\\ 3 & 4 & 2 & y\\ 5 & 6 & 3 &z\end{pmatrix} \sim  \begin{pmatrix} 1 & 0 & 0 & 0\\ 3 & 2 & 2 & y-3x\\ 5 & 4 & 4 & z-5x\end{pmatrix} \sim  \begin{pmatrix} 1 & 0 & 0 & 0\\ 3 & 2 & 0 & 0\\ 5 & 4 & x - 2y + z & 0\end{pmatrix}\sim  \begin{pmatrix} 1 & 0 & 0\\ 3 & 2 & 0\\ 5 & 4 & x - 2y + z\end{pmatrix}
   \]
   Alors l'équation cartésienne du sous-espace des colonnes est \(x - 2y + z = 0\).
   \item Cherchons le noyau et l'image de l'endomorphisme représenté par \(M\), on doit utiliser la mémorisation, ie:
   \[
      M = \begin{blockarray}{ccc} C_1 & C_2 & C_3 \\\begin{block}{(ccc)} 1 & 2 & 1\\3 & 4 & 2\\5 & 6 & 3\\\end{block}\end{blockarray} 
      \sim
      \begin{blockarray}{ccc} C_1 & C_2-2C_1 & 2C_3 - 2C_1 \\\begin{block}{(ccc)} 1 & 0 & 0\\3 & 2 & 2\\5 & 4 & 4\\\end{block}\end{blockarray} 
      \sim
      \begin{blockarray}{ccc} C_1 & C_2-2C_1 & 2C_3 - C_2 \\\begin{block}{(ccc)} 1 & 0 & 0\\3 & 2 & 0\\5 & 4 & 0\\\end{block}\end{blockarray} 
   \]
   On en conclut qu'une base de l'image est \([(1, 3, 5), (0, 2, 4)]\).\+
   Le noyau est de dimension 1 et par la mémorisation, une combinaison linéaire nulle est \(0C_1-C_2+2C_3\), le vecteur \((0, -1, 2)\) est bien dans le noyau et est donc une base du noyau.

\end{itemize}\<

\chapter*{\chapterstyle{II --- Déterminant}} % Fini 99%
\addcontentsline{toc}{section}{Déterminant}

Soit \(A = (a_{i,j})\) une matrice carrée de taille \(n\), nous cherchons à définir une application \(\phi: \mathcal{M}_n(\K) \longrightarrow \R\) telle que:
\[
   A \in \text{GL}_n(\K) \Longleftrightarrow \phi(A) = 0
\]
\subsection*{\subsecstyle{Définition{:}}}
\addcontentsline{toc}{subsection}{Définition}

On peut alors montrer\footnote[1]{Particulièrement difficile..} qu'une telle application existe qu'on appellera \textbf{déterminant}, qu'on note \(|A|\) et qu'on définit alors par récurrence:
\begin{itemize}
   \item Si \(n = 1\) alors \(|(a)| = a \) 
   \item Sinon \(|A| := a_{1, 1}|A_{1, 1}| - a_{1, 2}|A_{1, 2}| + \ldots + (-1)^{n+1}a_{1, n}|A_{1, n}|\)
\end{itemize}
Où la notation \(A_{1, 1}\) signifie la matrice \(A\) à laquelle on a retiré la première ligne et la première colonne. \<

Cette opération s'appelle aussi \textbf{développement selon la première ligne} du déterminant, elle se comprends visuellement par:
\[
   |A| = \begin{vmatrix} a & b & c \\ d & e & f \\ g & h & i \end{vmatrix} = 
   a \begin{vmatrix} \square & \square & \square \\ \square & e & f \\ \square & h & i \end{vmatrix} - 
   b \begin{vmatrix} \square & \square & \square \\ d & \square & f \\ g & \square & i \end{vmatrix} + 
   c \begin{vmatrix} \square & \square & \square \\ d & e & \square \\ g & h & \square \end{vmatrix}   
\]
On peut ainsi \textbf{développer le déterminant} selon n'importe quelle ligne ou colonne en suivant cette règle. On régresse ainsi vers des déterminants de plus petite taille, et aprés plusieurs étapes, vers un réel.

\subsection*{\subsecstyle{Propriétés{:}}}
\addcontentsline{toc}{subsection}{Propriétés}

On peut montrer que cette application est \textbf{une forme multilinéaire alternée} en les colonnes ou lignes de la matrice, en particulier, on a les propriétés suivantes:
\begin{itemize}
   \item Ajouter à une colonne une combinaison linéaire \textbf{des autres colonnes} ne change pas le déterminant.
   \item Echanger deux colonnes d'une matrice change le signe du determinant.
   \item Si deux colonnes sont égales ou qu'une des colonnes est nulle, le déterminant est nul.
\end{itemize}
En outre la multilinéarité sur les colonnes ou lignes implique:
\customBox{width=3cm}{
   \(|\lambda A| = \lambda^n |A|\)
}
Et on peut aussi montrer que le déterminant est multiplicatif, ie:
\customBox{width=4cm}{
   \(|A \times B| = |A|\times|B|\)
}

Si on considère une famille \(\Fam = (e_1 , \ldots, e_n)\) de vecteurs de coordonées \((C_1, \ldots, C_n)\) dans une certaine base \(\mathscr{B}\), alors on peut définir le déterminant \textbf{d'une famille de vecteurs} dans cette base par:
\[
   |\Fam|_{\mathscr{B}} = |(C_1, \ldots, C_n)|
\]
\begin{center}
   \textit{Le déterminant d'une famille par rapport à \(\mathscr{B}\) est alors simplement le déterminant de la matrice construite avec les coordonées de ses vecteurs dans la base \(\mathscr{B}\).}
\end{center}

Soit deux bases \(\mathscr{B, B'}\), on pourrait alors considèrer \textbf{l'effet d'un changement de base} sur un tel déterminant et on a alors:
\customBox{width=5cm}{
   \(
      |\Fam|_{\mathscr{B'}} = |\Fam|_{\mathscr{B}} \times |\mathscr{B'}|_{\mathscr{B}}
   \)
}

\subsection*{\subsecstyle{Cofacteurs{:}}}
\addcontentsline{toc}{subsection}{Cofacteurs}

Reprenons en simplifant la définition donnée plus haut du développement selon la première ligne:
\[
   |A| := a_{1, 1}|A_{1, 1}| - a_{1, 2}|A_{1, 2}| + \ldots + (-1)^{n+1}a_{1, n}|A_{1, n}| = \sum_{j=1}^{n}a_{1, j}(-1)^{j+1}|A_{1, j}|
\]
On appelle alors \textbf{cofacteur} de l'élément \(a_{i, j}\) le scalaire:
\customBox{width=5cm}{
   \(
      \text{cof}(a_{i, j}) := (-1)^{j+i}|A_{i, j}|
   \)
}   
Intuitivement en reprenant la visualisation exposée plus haut:
\[
   a \begin{vmatrix} \square & \square & \square \\ \square & \color{BrightRed1}e & \color{BrightRed1}f \\ \square & \color{BrightRed1}h & \color{BrightRed1}i \end{vmatrix} - 
   b \begin{vmatrix} \square & \square & \square \\ \color{BrightRed1}d & \square & \color{BrightRed1}f \\ \color{BrightRed1}g & \square & \color{BrightRed1}i \end{vmatrix} + 
   c \begin{vmatrix} \square & \square & \square \\ \color{BrightRed1}d & \color{BrightRed1}e & \square \\ \color{BrightRed1}g & \color{BrightRed1}h & \square \end{vmatrix}  
\]
Ce sont simplement les déterminants mineurs que l'on calcule à chaque itération \textbf{en tenant compte du signe qui précède le coefficient}.\+
En pratique pour trouver le signe du cofacteur, on ne calcule pas \((-1)^{i + j}\), mais on le détermine par \textbf{la règle de l'échiquier} qu'on peut se représenter comme suit:
\[
   \begin{vmatrix} + & - & + & - & \ldots \\ - & + & - & + &\ldots \\ + & - & + & - &\ldots \\ \vdots & \vdots & \vdots & \vdots \\\end{vmatrix}
\]

\underline{Exemple:} Pour \(\begin{pmatrix}
   1 & 2 & 3 \\ 4 & 5 & 6 \\ 7 & 8 & 9
\end{pmatrix}\), le cofacteur du coefficient en position \((1, 2)\) est 
\(- \begin{vmatrix} 4 & 6 \\ 7 & 9\end{vmatrix} \)

\subsection*{\subsecstyle{Comatrice{:}}}
\addcontentsline{toc}{subsection}{Comatrice}

On peut alors définir \textbf{la comatrice} d'une matrice \(A\) donnée, c'est en fait simplement \textbf{la matrice des cofacteurs} de \(A\), ie le terme en position \((i, j)\) de la comatrice est exactement le cofacteur de l'élement en position \((i, j)\). \<

L'intérêt principal de la comatrice est de calculer l'inverse de matrices inversibles, en effet, on peut montrer l'identité:
\[
   A \times \text{com}(A)^{\top} = |A|I_n   
\]
Et en particulier si le déterminant est non-nul (donc si \(A\) est inversible), on a:
\customBox{width=4.5cm}{
   \[
      A^{-1} = \frac{1}{|A|}\text{com}(A)^{\top}  
   \]
}   
Cette formule est néanmoins particulièrement inexploitable pour \(n > 3\) du fait de la quantité de calcul à réaliser.

\subsection*{\subsecstyle{Formules de Cramer{:}}}
\addcontentsline{toc}{subsection}{Formules de Cramer}

On considère un systême linéaire de \(n\) équations à \(n\) inconnues, alors on peut l'écrire matriciellement sous la forme:
\[
   AX = B   
\]
Pour \(A \in \mathcal{M}_n(\K)\), \(X\) un vecteur indéterminé, et \(B\) le vecteur fixé du second membre des équations.
\pagebreak

Un tel système est dit \textbf{système de Cramer} si la matrice \(A\) est inversible, il possède alors une unique solution qui s'écrit matriciellement \(X = A^{-1}B\). \<

Alors on montre que sa i-ème inconnue \(x_i\) (et par suite la solution tout entière X = \((x_1, \ldots, x_n)\)) est donnée par:
\customBox{width=2.5cm}{
   \[
      x_i = \frac{|A_i|}{|A|}
   \]
} 
Dans laquelle \(A_i\) désigne la matrice obtenue en remplaçant la i-ème colonne de \(A\) par le second membre \(B\)\<

\underline{Exemple:}
\[
   \begin{cases} 
      x + 4y + 8z = 2 \\
      2x + 5y + 8z = 2 \\
      3x + 6y + 9z = 3 \\
   \end{cases}   \Longleftrightarrow 
   \begin{pmatrix}
      1 & 4 & 8 \\
      2 & 5 & 8 \\
      3 & 6 & 9
   \end{pmatrix} 
   \begin{pmatrix}
      x \\
      y \\
      z  
   \end{pmatrix} =
   \begin{pmatrix}
      2 \\
      2 \\
      3  
   \end{pmatrix}
\]
Alors on a:
\[
   y = \frac{
      \begin{vmatrix}
         1 & \color{BrightRed1} 2 & 8 \\
         2 & \color{BrightRed1} 2 & 8  \\
         3 & \color{BrightRed1} 3 & 9
      \end{vmatrix}
   }{|A|} = \frac{6}{-3} = -2
\]

\subsection*{\subsecstyle{Interprétation géométrique{:}}}
\addcontentsline{toc}{subsection}{Interprétation géométrique}

Le déterminant admet un interprétation géométrique intéressante liée au concept \textbf{d'orientation} d'un espace et celui d'aire, de volumes, et d'hypervolumes de manière générale\footnote[1]{Voir le chapitre géométrie pour plus de détails, en particulier la partie sur le produit vectoriel.}. \<

Considérons le cas simple d'une base \(\mathscr{B} = (e_1, e_2)\) de vecteurs de \(\R^2\), alors le déterminant de cette base correspond \textbf{à l'aire algébrique\footnote[2]{C'est une aire ``signée'', ie l'aire géométrique est donc la valeur absolue de cette aire algébrique.} du parallélogramme formé par les deux vecteurs}, ie:
\begin{center}
   \begin{tikzpicture}[xscale=1.6, yscale=1.4]
      \draw[black!15] (-0.5,0) -- (4,0);
      \draw[black!15] (0, -0.5) -- (0,2);

      \draw[color=black!0, fill=BrightBlue1!25] (0,0) -- (0.6, 1.5) -- (3.1, 1.9) -- (2.5, 0.4);

      \draw node[color=DarkBlue1, rotate=8] at (1.55, 0.95) {$|\text{det}(e_1, e_2)|$};

      \draw[-latex, color=DarkBlue1, thick] (0, 0) -- (0.6, 1.5) node[left] {$e_1$};
      \draw[-latex, color=DarkBlue1, thick] (0, 0) -- (2.5, 0.4) node[below right] {$e_2$};

      \draw[color=DarkBlue1, dashed] (0.6, 1.5) -- (3.1, 1.9)-- (2.5, 0.4);


   \end{tikzpicture}
\end{center}
Plus généralement, le déterminant d'une famille de \(n\) vecteurs est \textbf{l'hypervolume algébrique} du parallélotope formé par ces \(n\) vecteurs.\<

Le lien avec l'orientation de l'espace quand à lui s'explique par le caractère \textbf{alterné} du déterminant.\+
En effet, en reprenant l'exemple ci-dessus, on voit que si on permute \(e_1\) et \(e_2\), alors l'aire géométrique reste la même, mais le déterminant change de signe, on dira alors que cette permutation a donc changé l'orientation de l'espace.\<

En particulier, on dira qu'une application linéaire \textbf{préserve l'orientation de l'espace}, si son déterminant est positif.

\chapter*{\chapterstyle{II --- Espaces affines}} % 90 Fini
\addcontentsline{toc}{section}{Espaces affines}
Soit \(\mathscr{E}\) un ensemble non-vide et \(V\) un \(\R\)-espace vectoriel de dimension \(n\) finie.
\subsection*{\subsecstyle{Définition {:}}}
On appelle \textbf{espace affine\footnote[1]{Ses éléments sont alors appelés des \textbf{points}.} de direction} \(V\) le couple \((\mathscr{E}, +)\) avec:
\[
   \begin{aligned}
      + : &&\mathscr{E} \times V &\longrightarrow \mathscr{E}\\
      &&(\,A, \, u\,) &\longmapsto A + u
   \end{aligned}
\]
La loi + doit vérifer les axiomes suivants 
\footnote[2]{Une action d'un groupe \((G, \star)\) sur \(E\) est une application \(+\) de \(G \times E \) dans \(E\) qui vérifie: \[
   \forall g, g' \in G  \; , \; \forall x \in E \; ; \; x + (g \star g') = (x + g) \star g'
\] \hspace{15pt} En d'autres termes, additionner d'abord les vecteurs, ou d'abord le point avec le vecteur n'importe pas.}:
\customBox{width=13cm}{
   \begin{align*}
      &\textbf{Existence d'un neutre} && \forall A \in \mathscr{E} \; ; \; A + 0_V = A\\
      &\textbf{Action de groupe} && \forall A \in \mathscr{E} \; , \; \forall u,v \in V \; ; \; A + (u + v) = (A + u) + v \\
      &\textbf{Unicité du translaté} && \forall A, B \in \mathscr{E} \; , \; \exists !u \in V \; ; \; A + u = B 
   \end{align*}
}    

Etant donné deux points \(A, B \in \mathscr{E}\), on note alors \(\overrightarrow{AB}\) l'unique vecteur \(u\) qui vérifie:
\customBox{width=2.5cm}{
   \(A + u = B\)
}

\subsection*{\subsecstyle{Propriétés {:}}}
Si \(\mathscr{E}\) est un espace affine de direction \(V\), on a alors les propriétés suivantes:
\customBox{width=8cm}{
   \begin{align*}
      &\textbf{Relation de Chasles} && \overrightarrow{AB} + \overrightarrow{BC} = \overrightarrow{AC} \\
      &\textbf{Existence d'un neutre} && \overrightarrow{AB} + \overrightarrow{BA} = O_V
    \end{align*}
}

\subsection*{\subsecstyle{Repères {:}}}
Un repère de \(\mathscr{E}\) est un couple \(\mathscr{R} = (O, \mathscr{B})\) formé d'un point \(O \in \mathscr{E}\) et d'une base \(\mathscr{B} = (e_1, \ldots, e_n)\) de \(V\). On appelle alors le point \(O\) \textbf{origine} du repère, et les vecteurs \(e_1, \ldots, e_n\) \textbf{vecteurs de base} du repère.\<

Soit \(i\in\inticc{1}{n}\), alors pour tout point \(A \in \mathscr{E}\), on appelle \textbf{coordonées} de \(A\) dans le repère \(\mathscr{R}\), les composantes \((x_i)_i\) du vecteur qui représente la translation de \(O\) vers \(A\) dans la base \(\mathscr{B}\), et on a la caractérisation élémentaire suivante:
\[
   \overrightarrow{OA} = x_1e_1 + \ldots + x_ne_n 
\]

\subsection*{\subsecstyle{Cas particulier des espaces vectoriels {:}}}
Il est important de noter que le couple \((V, +)\) forme un espace affine de \textbf{direction lui-même}. En effet, si on considère un élément de \(V\) comme un point, alors le couple  \((V, +)\) forme un espace affine de direction \(V\) et la loi \(+\) est alors exactement la loi de composition interne de \(V\).\<

L'unique vecteur \(u\) tel que \(A + u = B\) est exactement \(B - A\) (ici \(A, B\) sont des points de \(V\) qui se trouvent être des vecteurs dans ce cas particulier). \<\<

\chapter*{\chapterstyle{II --- Espaces euclidiens}} % 90 Fini
\addcontentsline{toc}{section}{Espaces euclidiens}

Soit \(E\) un \(\R\)-espace vectoriel de dimension \(n\) finie. On considère \(E\) avec sa structure d'espace affine canonique.
On appelle \textbf{espace euclidien} un tel espace vectoriel si il est muni d'un \textbf{produit scalaire}.\<

Soit \(u, v\) deux vecteurs de \(E\) et \((u_i)_{i \leq n}\) et \((v_i)_{i \leq n}\) les composantes de \(u\) et de \(v\) dans la base canonique de \(E\).

\subsection*{\subsecstyle{Vecteurs {:}}}
Soit deux points \(A, B\), on note \(\overrightarrow{AB}\) l'unique vecteur \(u\) de \(E\) tel que \(A + u = B\), on appelle alors \(A\) le \textbf{point initial} et \(B\) le \textbf{point final}. \<

On peut alors considérer l'application \(\phi\) qui associe aux points \(A, B\) le vecteur \(\overrightarrow{AB}\) en \textbf{position standard}, et on a:
\[
   \begin{aligned}
      \phi: E \times E &\longrightarrow E\\
      (A, B) &\longmapsto B - A
   \end{aligned}
\]
Ici, \(A, B\) sont des points qui trouvent aussi être des vecteurs, car l'espace affine sous-jacent est un espace vectoriel, et donc la soustractrion ci-dessus est bien définie.

\subsection*{\subsecstyle{Produit scalaire {:}}}
Soit \(\phi : E^2 \longrightarrow \R\) une forme \textbf{bilinéaire}. On dit que l'application \(\phi\) est un \textbf{produit scalaire} si et seulement si il vérifie:

\customBox{width=7cm}{
   \begin{align*}
      &\textbf{Symétrique} &&\dotproduct{u}{v} = \dotproduct{v}{u} \\
      &\textbf{Définie} &&\dotproduct{u}{u} = 0 \implies u = 0 \\
      &\textbf{Positive} &&\dotproduct{u}{v} \geq 0
   \end{align*}
}   


Par exemple, si on considère \(E = \R^n\), alors on peut définir le produit scalaire canonique de \(\R^n\) comme ci-dessous:
\[
   \phi : (u, v) \longmapsto \sum_{k=1}^{n} u_iv_i 
\]
Par ailleurs, on a \textbf{l'inégalité de Cauchy-Schwartz}:
\[
    |\dotproduct{u}{v}| \leq \vectNorm{u} \vectNorm{v}
\]

\subsection*{\subsecstyle{Norme{:}}}
On appelle \textbf{norme} une application \(\vectNorm{\cdot}\) de E vers \(\R\) qui vérifie:
\customBox{width=8.5cm}{
   \begin{align*}
      &\textbf{Séparation} &&\vectNorm{u} = 0 \implies u = 0\\
      &\textbf{Absolue Homogénéité} &&\vectNorm{\lambda u} = |\lambda| \vectNorm{u}\\
      &\textbf{Sous-Additivité} &&\vectNorm{u + v} \leq \vectNorm{u} + \vectNorm{v}
   \end{align*}
}    

Un espace vectoriel muni d'une norme est alors appelé \textbf{espace vectoriel normé}. Dans le cas d'un espace euclidien, on définit la norme euclidienne usuelle:
\[
   \vectNorm{u} = \sqrt{\sum_{i=1}^n u_i^2}
\]
Cette norme induit alors une \textbf{métrique} sur \(E\)  comme suit:
\[
   \begin{aligned}
      d: E \times E &\longrightarrow E\\
      (u, v) &\longmapsto \vectNorm{v - u}
   \end{aligned}
\]
Muni de cette métrique, l'espace euclidien \(E\) est alors un \textbf{espace métrique homogène}.

\subsection*{\subsecstyle{Angles{:}}}
Dans un espace euclidien, on définit \textbf{l'angle} \(\theta := \angle AOC\) comme étant l'angle entre les bipoints \(OA\) et \(OC\) et on a alors la formule de \textbf{Pythagore généralisée}:
\[
   \vectNorm{BC}^2 = \vectNorm{AB}^2 + \vectNorm{AC}^2 - 2\vectNorm{AB}\vectNorm{AC}\cos(\theta)
\]
Puis par suite on peut montrer la caractérisation du produit scalaire:
\customBox{width=4.6cm}{
   \(\dotproduct{u}{v} = \vectNorm{u} \vectNorm{v} \cos{(\theta)}\)
}

\subsection*{\subsecstyle{Composantes scalaires et projections{:}}}
Soit \(\alpha\) l'angle entre les deux vecteurs \(u\) et \(v\).
On appelle \textbf{composante scalaire} de \(u\) par rapport à \(v\) le scalaire:
\[
   \text{comp}_v(u) = \vectNorm{u}\cos(\alpha)   
\]
On définit alors la \textbf{projection orthogonale} de \(u\) sur \(v\):
\customBox{width=5.1cm}{
   \[\text{proj}_v(u) = \text{comp}_v(u)\frac{v}{\vectNorm{v}}\]
}
Visuellement, si on considère deux vecteurs de \(E = \R^2\), on a:
\begin{center}
   \begin{tikzpicture}[
      >=stealth,scale=1,line cap=round,
      bullet/.style={circle,inner sep=0.5pt,fill}
   ]
      \draw[->] (-1,0) -- (9,0);
      \draw[->] (0,-1) -- (0, 5.5);
   
      \draw foreach \X in {3,6}{
         (\X, 0.1) -- (\X, -0.1) node[below]{$\X$}
      };
      \draw foreach \Y in {2,4}{
         (0.1,\Y) -- (-0.1,\Y) node[left]{$\Y$}
      };
   
      \draw[->]  (0,0) coordinate (O) -- (3,4) coordinate (B) node[above]{$u$};
      \draw[->]  (0,0) coordinate (O) -- (6,2) coordinate (A) node[above]{$v$};

      \draw[dashed, color=DarkBlue1] (3, 4) -- (3.9, 1.3);

      \draw[->, line width=2pt, color=DarkBlue1]  (0,0) -- (3.9, 1.3);
      \path (2,1.2) node[color=DarkBlue1]{$\text{proj}_v(u)$};
      \draw[solid, line width=2pt, color=BrightBlue1]  (0,0) -- (4.111, 0) node[above]{$\text{comp}_v(u)$};
      \draw pic["$\alpha$", draw, gray, angle radius=1cm] {angle = A--O--B};

   \end{tikzpicture}   
\end{center}
On remarque alors que \textbf{la composante scalaire de \(u\) est exactement la norme du projeté de \(u\) sur \(v\)}

\subsection*{\subsecstyle{Orientation{:}}}
Pour la suite du chaptire, on se place dans \(E = \R^3\), soit deux vecteur \(u, v\) non parallèles, on appelle \textbf{orientation directe} l'orientation de l'espace suivante \textbf{la règle de la main droite}, et orientation indirecte celle suivant la règle de la main gauche.\<

Plus formellement, l'orientation d'un espace euclidien de base \(\Fam\) est donné par le signe du déterminant de cette base:
\begin{itemize}
   \item Si le déterminant est positif, la base est en \textbf{orientation directe}.
   \item Si le déterminant est négatif, la base est en \textbf{orientation indirecte}.
\end{itemize}
\pagebreak

\subsection*{\subsecstyle{Produit vectoriel{:}}}
On définit le \textbf{produit vectoriel} de deux vecteurs \(u, v\) comme l'unique vecteur \(w\) de \(\R^3\) qui vérifie:
\customBox{width=6.5cm}{
    \begin{center}
      $\vectNorm{w} = \vectNorm{u}\vectNorm{v}|\sin(\theta)|$\\
      $w$ \textbf{est orthogonal à }$u$\textbf{ et à }$v$\\
      \textbf{La base } $(u, v, w)$ \textbf{ est directe}
    \end{center}
}

C'est une application de \(E \times E \longrightarrow E\) qui vérfie:
\customBox{width=9.5cm}{
   \begin{align*}
      &\textbf{Homogénéité} && \lambda_1 u \times v = \lambda_1 (u \times v)\\
      &\textbf{Anticommutativité} && u \times v = - (v \times u)\\
      &\textbf{Distributivité} && u \times (v + w) = u \times v + u \times w
    \end{align*}
}

Néanmoins, le produit vectoriel n'est en général \textbf{pas associatif}. En particulier, on a:
\[
    u \times (v \times w) = \dotproduct{u}{w}v - \dotproduct{u}{v}w
\]

\underline{Exemple:} On note \(\big|M\big|\) le déterminant d'une matrice, alors on a une caractérisation utlie pour calculer des produits vectoriels:
\[
   \begin{cases}
      u := (u_1, u_2, u_3)\\ 
      v := (v_1, v_2, v_3)
   \end{cases} \implies 
   u \times v = \Biggl(\begin{vmatrix}u_2 & u_3 \\ v_2 & v_3\end{vmatrix} ; -\begin{vmatrix}u_1 & u_3 \\ v_1 & v_3\end{vmatrix} ; \begin{vmatrix}u_1 & u_2 \\ v_1 & v_2\end{vmatrix}\Biggl)
\]

\subsection*{\subsecstyle{Volumes \& Produit triple{:}}}
On peut remarque que la norme du vecteur resultant d'un produit vectoriel est \textbf{l'aire du parallélogramme} formé par \(u\) et \(v\).\<

On peut alors montrer que le produit \(\dotproduct{u \times v}{w}\) est \textbf{le volume du parallélépipède} formé par \(u, v, w\). En effet, on a:
\[
   \dotproduct{u \times v}{w} = \vectNorm{u \times v} \vectNorm{w} \cos(\alpha) = \vectNorm{u \times v}h
\]
Avec \(\vectNorm{u \times v}\) qui est l'aire de la base est \(h\) la hauteur du parallélépipède.

\begin{center}
   \begin{tikzpicture}[
      >=stealth,scale=1,line cap=round,
      bullet/.style={circle,inner sep=0.5pt,fill}
   ]
      \coordinate (A) at (0, 0, 0);
      \coordinate (B) at (8, 0, 0);
      \coordinate (C) at (8, 0, -8);
      \coordinate (D) at (0, 0, -8);
      \coordinate (E) at (0, 2, -10);
      \coordinate (F) at (0, 2, -2);
      \coordinate (G) at (8, 2, -2);
      \coordinate (H) at (8, 2, -10);

      \draw[->] (A) -- (B) node[below right]{\(u\)};
      \draw[->] (A) -- (D) node[above left]{\(v\)};
      \draw[->] (A) -- (F) node[below right]{\(w\)};

      \draw[dashed, color=gray] (B) -- (C) -- (D) -- (E) -- (F) -- (G) -- (H) -- (E);
      \draw[dashed, color=gray] (B) -- (G);
      \draw[dashed, color=gray] (C) -- (H);
      \draw[dashed, color=black!10, fill=BrightBlue1, opacity=0.2] (E) -- (F) -- (G) -- (H) -- cycle;
      \draw[dashed, color=black!10, fill=BrightBlue1, opacity=0.15] (A) -- (B) -- (C) -- (D) -- cycle;
      \draw[dashed, color=black!10, fill=BrightBlue1, opacity=0.2] (A) -- (B) -- (G) -- (F) -- cycle;
      \draw[dashed, color=black!10, fill=BrightBlue1, opacity=0.2] (B) -- (C) -- (H) -- (G) -- cycle;

      \draw[->] (A) -- (0, 4, 0) coordinate (X) node[below right]{\(u \times v\)};
      \draw[gray] pic ["$\alpha$", draw, angle radius=2cm] {angle = F--A--X};

      \draw[->, color = DarkBlue3, line width = 1.2pt] (0, 0.01, 0) -- (0, 2.7, 0) coordinate (X) node[left]{\(h\)};
   \end{tikzpicture}   
\end{center}
Un autre manière de déterminer le volume d'un parallélépipède est de calculer le déterminant \(\begin{vmatrix}u_1 & v_1 & w_1 \\ u_2 & v_2 & w_2 \\ u_3 & v_3 & w_3\end{vmatrix}\)

\chapter*{\chapterstyle{II --- Droites \& Plans}} % 90 Fini
\addcontentsline{toc}{section}{Droites \& Plans}

Dans l'espace vectoriel \(\R^3\), nous savons qu'on peut définir des droites et des plans vectoriels, dans ce chapitre nous allons examiner comment exprimer mathématiquement les droites et plans plus généralement \textbf{affines} et calculer des distances entre ces figures ou des intersections. \<

Dans la suite, on considère un point \(A\) de coordonées \((x_0, y_0, z_0)\) et une vecteur \(\vv{v}\) de composantes \((a, b, c)\).

\subsection*{\subsecstyle{Droites{:}}}
Une droite de l'espace en trois dimensions est caractérisée par la donnée d'un \textbf{point} sur la droite, et d'un \textbf{vecteur directeur} de cette droite.\<

Si on considère une fonction vectorielle \(f(t)\) telle que \(f(0)\) soit la position de \(A\), on peut alors montrer que les points de la droite sont donnés par:
\customBox{width=3.5cm}{
   \(f(t) = f(0) + t\vv{v}\)
}
Cette fonction vectorielle est un cas particulier de \textbf{courbe paramétrée}, et cette notion est formalisée dans le chapitre au nom correspondant.\footnote[1]{On peut développer cette expression en composantes pour avoir une forme plus explicite.\label{developpement}} 

\subsection*{\subsecstyle{Plans{:}}}
Un plan de l'espace en trois dimensions peut être caractérisé par la donnée d'un \textbf{point} sur le plan et de \textbf{deux vecteurs} engendrant le plan.\<

Si on considère une fonction vectorielle \(f(t)\) telle que \(f(0)\) soit la position de \(A\), on peut alors montrer que les points de la droite sont donnés par:
\customBox{width=5cm}{
   \(f(t, t') = f(0) + t\vv{u} + t'\vv{v}\)
}
Cette fonction vectorielle est un cas particulier de \textbf{surface paramétrée}, et cette notion est formalisée dans le chapitre au nom correspondant.\footref{developpement}\<

Mais un plan de l'espace en trois dimensions peut aussi être caractérisé par la donnée d'un \textbf{point} sur le plan, et d'un \textbf{vecteur normal} au plan, en effet un plan est l'ensemble des vecteurs orthogonaux à un vecteur normal donné, on en déduit que les points \(P\) du plan \(\mathscr{P}\) vérifient l'équation vectorielle:
\customBox{width=3.5cm}{
   \(\dotproduct{\vv{v}}{\vv{AP}} = 0  \)
}

Si on déploie le produit sclaire en composantes, on obtient que \textbf{les coordonées} des points du plan vérifient:
\[
   a(x - x_0) + b(y - y_0) + c(z - z_0) = 0 \Longleftrightarrow ax + by + cz = C
\]

\begin{center}
   \textit{
      On peut alors remarque que les plans et les droites de l'espace affines sont exactement des translations des droites et plans linéaires de l'espace vectoriel sous-jacent. 
   }
\end{center}

\pagebreak

\subsection*{\subsecstyle{Distances \& Angles{:}}}
Nous sommes amenés à calculer des distances et des angles en géométrie euclidienne, par exemple calculer la distance entre un point \(P\) et une droite \(\mathscr{D}\).
Si nous connaissons un point \(A\) de la droite et un vecteur directeur \(\vv{v}\) de la droite, il suffit alors de calculer:
\begin{flalign*}
   &&\text{\customBox{width = 3cm}{$\vectNorm{\vv{AP}}\sin(\alpha)$}} &\hspace{-1cm}\shorteqnote{\raisebox{0.35cm}{\(\color{DarkBlue1}\bigstar\)}}
\end{flalign*}


L'apparition du sinus s'explique par le fait que la distance entre \(P\) et \(\mathscr{D}\) est \textbf{invariante par rotation}, et alors on peut se ramener à un problème trigonométrique sur un cercle de rayon \(\vectNorm{\vv{AP}}\).\+
Visuellement, on a:
\begin{center}
   \begin{tikzpicture}[
      >=stealth,scale=1,line cap=round,
      bullet/.style={circle,inner sep=0.5pt,fill}
   ]
      \coordinate (O) at (0,0);
      \coordinate (A) at (6,2);
      \coordinate (A') at (4.5, 1.5);
      \coordinate (P) at (3,4);

      \coordinate (V') at (-1.5, 4.5);
      \coordinate (Proj1) at (3.9, 1.3);
      \coordinate (Proj2) at (-0.9, 2.7);

      \draw[]  (-2, -2/3)  -- (8, 8/3) coordinate (A) node[above]{$\mathscr{D}$};

      \draw[->, color=DarkBlue1, line width = 1pt]  (O)  node[below]{$A$} -- (P) node[above]{$P$};
      \draw[->, color=DarkBlue1, line width = 1pt]  (O) -- (A') node[above]{$v$};
      \draw[->, color=DarkBlue1, line width = 1pt]  (O) -- (V') node[above]{$v^\perp$};
      \draw[solid, line width=2pt, color=BrightBlue1] (P) -- (Proj1);

      \draw node[color=BrightBlue1] at (-1.8, 1.7) {$\vectNorm{\vv{AP}}\sin(\alpha)$};

      \draw[solid, line width=2pt, color=BrightBlue1]  (O) -- (Proj2);
      \draw[dashed, color=DarkBlue1] (Proj2) -- (P);

      \draw pic["$\alpha$", draw, gray, angle radius=1cm] {angle = A--O--P};
      \draw pic["$\alpha'$", draw, gray, angle radius=1.2cm] {angle = P--O--Proj2};

      \RightAngle{(-2, -2/3)}{(0,0)}{(-1.5, 4.5)};
      \tkzMarkSegment[color=black, pos=0.5, size=5pt, mark=||](O, Proj2)
      \tkzMarkSegment[color=black, pos=0.5, size=5pt,mark=||](P, Proj1)

   \end{tikzpicture}   
\end{center}
On remarque alors une symétrie conceptuelle avec la projection vue plus haut: 
\begin{center}
   \textit{
      Chercher la distance entre un point \(P\) et une droite contenant \(A\) revient à\+ projeter \(\vv{AP}\) sur un vecteur orthogonal à la droite (en utilisant l'angle complémentaire).
   }
\end{center}
Par ailleurs, de la formule \(\color{DarkBlue1}(\bigstar)\) et de la définition du produit vectoriel, on peut en déduire une autre expression de la distance:
\customBox{width=3cm}{
   \[
      \frac{\vectNorm{\vv{AP} \times \vv{v}}}{\vectNorm{\vv{v}}}  
   \]
}
Sachant l'aire du parallélogramme et un de ses cotés, cette formule consiste exactement à retrouver la hauteur du parallèlogramme, qui constitura ici la distance recherchée. \<

Nous pouvons aussi être amenés à calculer la distance entre un point \(P\) et un plan \(\mathscr{P}\)\+
Si on connait un point \(A\) et un vecteur normal \(\vv{v}\) du plan, alors on peut trouver l'angle entre \(\vv{AP}\) et \(\vv{v}\) et il suffit de calculer la composante scalaire du vecteur \(\vv{AP}\) par rapport au vecteur normal.\+
Visullement cela donne:
\begin{center}
   \begin{tikzpicture}[
      >=stealth,scale=1,line cap=round,
      bullet/.style={circle,inner sep=0.5pt,fill}
   ]
      \coordinate (P1) at (0, 0, 0);
      \coordinate (P2) at (12, 0, 0);
      \coordinate (P3) at (12, 0, -8);  
      \coordinate (P4) at (0, 0, -8);      
    
      \coordinate (X) at (3, 1, -4);

      \coordinate (V1) at (8, 0, -4);
      \coordinate (V2) at (8, 2, -4);

      \coordinate (Proj1) at (8, 1, -4);
      \coordinate (Proj2) at (3, 0, -4);

      \draw[->] (V1) node[below right]{\(A\)} -- (V2) node[above left]{\(\vv{v}\)};
      \draw[->] (V1) -- (X) node[above left]{\(P\)};

      \draw[solid, line width=1pt, color=BrightBlue1] (V1) -- (Proj1);
      \draw node[color=BrightBlue1] at (9.25, 0.5, -4) {\(\vectNorm{\vv{AP}}\cos(\alpha)\)};
      \draw[dashed, color=gray] (X) -- (Proj1);
      \draw[solid, line width=1pt, color=BrightBlue1] (X) -- (Proj2);
      \draw[dashed, color=gray] (Proj2) -- (V1);

      %4,58
      \draw[->] (V1) -- (X) node[above left]{\(P\)};

      \draw[dashed, color=gray] (P1) -- (P2) -- (P3) -- (P4) -- (P1);
      \draw[dashed, color=black!10, fill=BrightBlue1, opacity=0.15] (P1) -- (P2) -- (P3) -- (P4) -- cycle;
      \RightAngle{(10, 0, -4)}{(V1)}{(V2)};
      \RightAngle{(X)}{(Proj2)}{(V1)};

      \draw pic["$\alpha$", draw, gray, angle radius=0.8cm] {angle = V2--V1--X};
      \tkzMarkSegment[color=black, pos=0.5, size=3pt, mark=||](Proj2, X)
      \tkzMarkSegment[color=black, pos=0.5, size=3pt, mark=||](V1, Proj1)

   \end{tikzpicture}   
\end{center}\<

Finalement, il peut être utile de savoir calculer \textbf{l'angle entre deux plans}, il suffit alors de calculer \textbf{l'angle entre leurs vecteurs normaux}.

